\chapter{Methodology}
\label{chapter:methodology}

This chapter introduces the mathematical construction used to turn a mean-field control problem into a model-free policy-gradient objective. We first fix the canonical controlled process and its policy parametrization, then define the randomized population perturbation whose limiting behavior justifies the transport estimators developed in the next chapter.

To keep the chapter focused on the construction and on the consequences of the main results, the more technical proofs are collected in Appendix~\ref{appendix:technical-proofs}. The main text keeps the statements and proof ideas needed to follow how the method fits together.

\section{Canonical construction of the control problem}

Fix a horizon $T\in\N^\star$ and denote by $\Tc=\{0,1,\ldots,T\}$ the set of time indices. Let $\Xc$ be the state space and let $A$ be the one-step action space, both assumed to be Polish. Given a random variable $Y$ on a probability space $(\Omega,\Fc,\P)$, we denote by $\P_Y$ or $\Lc(Y)$ its law under $\P$.

The notation $A$ is reserved for individual actions, whereas $\Ac$ denotes the set of admissible control processes. Given $\alpha=(\alpha_t)_{t=0}^{T-1}\in\Ac$ and an initial condition $\xi\sim\mu_0$, the controlled state process satisfies
\begin{align}
    \begin{cases}
        X_0=\xi,\\
        X_{t+1}=F_t(X_t,\P_{X_t},\alpha_t,\epsilon_{t+1}),
    \end{cases}
    \qquad t=0,\ldots,T-1,
\end{align}
or equivalently $\Lc\left(X_{t+1}\mid X_t,\P_{X_t},\alpha_t\right)=P_t(\cdot\mid X_t,\P_{X_t},\alpha_t)$. The reward functional is
\begin{align}
    \label{eq:methodology-control_problem}
    \sup_{\alpha\in\Ac}
    J(\alpha)
    \coloneqq
    \E\left[
        \sum_{t=0}^{T-1} r_t(X_t,\P_{X_t},\alpha_t)
        +g(X_T,\P_{X_T})
    \right].
\end{align}

In the model-free setting, controls are sampled from randomized Markov policies. We denote by $\Pi=\left\{\pi=(\pi_t)_{t=0}^{T-1}:\pi_t:\Xc\times\Pc(\Xc)\to\Pc(A)\text{ is measurable}\right\}$ the set of admissible policies, and define $\Ac\coloneqq\left\{\alpha=(\alpha_t)_{t=0}^{T-1}:\Lc(\alpha_t\mid X_t,\P_{X_t})=\pi_t(\cdot\mid X_t,\P_{X_t})\text{ for some }\pi\in\Pi\right\}$. The control problem \eqref{eq:methodology-control_problem} can then be written directly over policies:
\begin{align}
    \label{eq:methodology-equivalent_control}
    \sup_{\pi\in\Pi}
    J(\pi)
    \coloneqq
    \E^\pi\left[
        \sum_{t=0}^{T-1}r_t(X_t,\P_{X_t},\alpha_t)
        +g(X_T,\P_{X_T})
    \right],
\end{align}
where $\E^\pi$ denotes expectation under the trajectory law induced by the policy $\pi$.

To solve \eqref{eq:methodology-equivalent_control} by policy gradient methods, we restrict to a parametrized family $\theta\mapsto\pi^\theta=(\pi_t^\theta)_{t=0}^{T-1}$, with $\theta\in\Theta\subset\R^k$. We write $X^\theta$ and $\alpha^\theta$ for the corresponding state and action processes, and the finite-dimensional objective becomes
\begin{align}
    \sup_{\theta\in\Theta}J(\theta)
    =
    J(\pi^\theta).
\end{align}

The difficulty is that the policy and the reward depend on the population law $(\P_{X_t^\theta})_{t\in\Tc}$, whose sensitivity with respect to $\theta$ is not directly available in a model-free setting. The construction below introduces an exogenous randomization of these laws. The perturbed objective will converge to the original one as the perturbation scale $\lambda$ tends to zero, while providing a likelihood-ratio handle on the missing population derivative.

We work on a probability space rich enough to support the following independent random variables:
\begin{enumerate}
    \item an initial variable $\xi\sim\mu_0$;
    \item transition noises $(\epsilon_t)_{t=1}^{T}$ with common law $\rho\in\Pc(E)$;
    \item policy noises $(\tilde\epsilon_t)_{t=0}^{T-1}$ with common law $\tilde\rho\in\Pc(\tilde E)$;
    \item perturbation variables $(f_t)_{t=0}^{T}$ with common law $\nu\in\Pc(\mathsf F)$.
\end{enumerate}
By the representation theorem for probability kernels of Kallenberg \cite{kallenberg1997foundations}, the transition kernel can be represented by measurable maps $\Xc\times\Pc(\Xc)\times A\times E \ni(x,m,a,e)\mapsto F_t(x,m,a,e)$ such that $F_t(x,m,a,\cdot)\sharp\rho=P_t(\cdot\mid x,m,a)$. Similarly, each randomized policy can be represented by measurable maps $\Xc\times\Pc(\Xc)\times\tilde E\ni(x,m,\tilde e)\mapsto\tilde F_t^\theta(x,m,\tilde e)$ such that $\tilde F_t^\theta(x,m,\cdot)\sharp \tilde\rho= \pi_t^\theta(\cdot\mid x,m)$.

\begin{remark}[Canonical construction of the problem]
    The preceding variables can be realized on the canonical product space $\Omega=\Xc\times E^T\times\tilde E^T\times\mathsf F^{T+1}$, endowed with its Borel $\sigma$-algebra and with the product measure
    \begin{align}
        \P(\d\omega)
        =
        \mu_0(\d x_0)
        \rho^{\otimes T}(\d\epsilon_{1:T})
        \tilde\rho^{\otimes T}(\d\tilde\epsilon_{0:T-1})
        \nu^{\otimes(T+1)}(\d f_{0:T}).
    \end{align}
\end{remark}

For a fixed policy $\pi^\theta$, the unperturbed canonical process is defined by
\begin{align}
    \begin{cases}
        X_0^\theta(\omega)=x_0,\\
        \alpha_t^\theta(\omega)
        =
        \tilde F_t^\theta
        \left(X_t^\theta(\omega),\P_{X_t^\theta},\tilde\epsilon_t(\omega)\right),\\
        X_{t+1}^\theta(\omega)
        =
        F_t
        \left(X_t^\theta(\omega),\P_{X_t^\theta},\alpha_t^\theta(\omega),\epsilon_{t+1}(\omega)\right),
    \end{cases}
    \qquad t=0,\ldots,T-1.
\end{align}
By construction, it follows that
\begin{align}
    \begin{cases}
        X_0^\theta\sim\mu_0,\\
        \Lc\left(\alpha_t^\theta\mid X_t^\theta,\P_{X_t^\theta}\right)
        =
        \pi_t^\theta(\cdot\mid X_t^\theta,\P_{X_t^\theta}),\\
        \Lc\left(X_{t+1}^\theta\mid X_t^\theta,\P_{X_t^\theta},\alpha_t^\theta\right)
        =
        P_t(\cdot\mid X_t^\theta,\P_{X_t^\theta},\alpha_t^\theta).
    \end{cases}
\end{align}

\section{Perturbed process and value function}

We now randomize the population argument while keeping the reference population flow $(\P_{X_t^\theta})_{t\in\Tc}$. For the abstract construction, assume that each perturbation variable $f\in\mathsf F$ induces a measurable map $T_f^\lambda:\Xc\to\Xc$. In the vector-space case used by the continuous benchmarks, this is written as $T_f^\lambda=I_d+\lambda f$. For any $\theta\in\Theta$ and $\lambda\geq0$, define
\begin{align}
    \Phi_t^{\lambda,\theta}(f)
    \coloneqq
    T_f^\lambda\sharp\P_{X_t^\theta},
    \qquad
    \P_t^{\lambda,\theta}(\omega)
    \coloneqq
    \Phi_t^{\lambda,\theta}(f_t(\omega)).
\end{align}
The perturbed trajectory associated with $\theta$ is then
\begin{align}
    \begin{cases}
        X_0^{\lambda,\theta}(\omega)=x_0,\\
        \alpha_t^{\lambda,\theta}(\omega)
        =
        \tilde F_t^\theta
        \left(X_t^{\lambda,\theta}(\omega),\P_t^{\lambda,\theta}(\omega),\tilde\epsilon_t(\omega)\right),\\
        X_{t+1}^{\lambda,\theta}(\omega)
        =
        F_t
        \left(X_t^{\lambda,\theta}(\omega),\P_t^{\lambda,\theta}(\omega),\alpha_t^{\lambda,\theta}(\omega),\epsilon_{t+1}(\omega)\right),
    \end{cases}
    \qquad t=0,\ldots,T-1.
\end{align}
Conditionally on the realized perturbation, this process satisfies
\begin{align}
    \label{eq:methodology-perturbed_equations}
    \begin{cases}
        X_0^{\lambda,\theta}\sim\mu_0,\\
        \Lc\left(\alpha_t^{\lambda,\theta}
        \mid X_t^{\lambda,\theta},\P_t^{\lambda,\theta}\right)
        =
        \pi_t^\theta(\cdot\mid X_t^{\lambda,\theta},\P_t^{\lambda,\theta}),\\
        \Lc\left(X_{t+1}^{\lambda,\theta}
        \mid X_t^{\lambda,\theta},\P_t^{\lambda,\theta},\alpha_t^{\lambda,\theta}\right)
        =
        P_t(\cdot\mid X_t^{\lambda,\theta},\P_t^{\lambda,\theta},\alpha_t^{\lambda,\theta}).
    \end{cases}
\end{align}
The perturbed reward functional is
\begin{align}
    J^\lambda(\theta)
    \coloneqq
    \E\left[
        \sum_{t=0}^{T-1}
        r_t(X_t^{\lambda,\theta},\P_t^{\lambda,\theta},\alpha_t^{\lambda,\theta})
        +
        g(X_T^{\lambda,\theta},\P_T^{\lambda,\theta})
    \right],
    \qquad
    \sup_{\theta\in\Theta}J^\lambda(\theta).
\end{align}

Equivalently, write
\begin{align}
    s
    &=
    (x_0,f_0,a_0,\ldots,x_{T-1},f_{T-1},a_{T-1},x_T,f_T),\\
    \tilde s
    &=
    (x_0,m_0,a_0,\ldots,x_{T-1},m_{T-1},a_{T-1},x_T,m_T).
\end{align}
The trajectory law can be expressed either in terms of the perturbation variables $f_t$,
\begin{align}
    p_\theta^\lambda(\d s)
    &=
    \mu_0(\d x_0)
    \prod_{\tau=0}^{T-1}
    \nu(\d f_\tau)
    \pi_\tau^\theta(\d a_\tau\mid x_\tau,\Phi_\tau^{\lambda,\theta}(f_\tau))
    P_\tau(\d x_{\tau+1}\mid x_\tau,\Phi_\tau^{\lambda,\theta}(f_\tau),a_\tau)
    \cdot\nu(\d f_T),
\end{align}
or in terms of the induced random population arguments $m_t=\Phi_t^{\lambda,\theta}(f_t)$,
\begin{align}
    \tilde p_\theta^\lambda(\d\tilde s)
    &=
    \mu_0(\d x_0)
    \prod_{\tau=0}^{T-1}
    \left(\Phi_\tau^{\lambda,\theta}\sharp\nu\right)(\d m_\tau)
    \pi_\tau^\theta(\d a_\tau\mid x_\tau,m_\tau)
    P_\tau(\d x_{\tau+1}\mid x_\tau,m_\tau,a_\tau)
    \cdot
    \left(\Phi_T^{\lambda,\theta}\sharp\nu\right)(\d m_T).
\end{align}
Therefore
\begin{align}
    J^{\lambda}(\theta)
    &=
    \int
    \left[
        \sum_{t=0}^{T-1}
        r_t(x_t,\Phi_t^{\lambda,\theta}(f_t),a_t)
        +
        g(x_T,\Phi_T^{\lambda,\theta}(f_T))
    \right]
    p_{\theta}^\lambda(\d s)\\
    &=
    \int
    \left[
        \sum_{t=0}^{T-1} r_t(x_t,m_t,a_t)
        +
        g(x_T,m_T)
    \right]
    \tilde p_{\theta}^\lambda(\d\tilde s).
\end{align}

\section{Discrete-state space}

Let $\Xc=\{1,\ldots,N\}$. Define the probability simplex on $\Xc$ as $\Delta_N\coloneqq\{p\in\R_+^N : \sum_{i=1}^p p_i=1\}$ and denote $\Delta_N^\circ$ its relative interior. For a policy parameter $\theta\in\Theta$, let $\mu_t^\theta(i)\coloneqq \P(X_t^\theta=i)$ be the law of the unperturbed controlled state process, for $i\in\{1,\ldots,N\}$.

We introduce a random perturbation of the population law $\mu_t^\theta$ while keeping the perturbed object inside the simplex. Consider a stochastic matrix $Q_t = (Q_t^{i,j})_{1\leq i,j\leq N}$, and focus on the rank-one perturbation $Q_t^{j,i} = q_t^i$ where $q_t$ is a random element of $\Delta_N^\circ$. The perturbed population law is then
\begin{align}
    M_t^{\lambda,\theta}(\omega)\coloneqq (1-\lambda)\mu_t^\theta + \lambda q_t(\omega).
\end{align}
Notice that $M_t^{\lambda,\theta}$ is a random probability vector, but it should not be identified with the law of the perturbed state process defined below.

We now define the perturbed controlled process
\begin{align}
    \begin{cases}
        X_0^{\lambda,\theta} &\sim \mu_0, \\
        \alpha_t^{\lambda,\theta} &\sim \pi_t^\theta(\cdot\mid X_t^{\lambda,\theta}, M_t^{\lambda,\theta}), \\
        X_{t+1}^{\lambda,\theta} &\sim P_t(\cdot\mid X_t^{\lambda,\theta}, M_t^{\lambda,\theta}, \alpha_t^{\lambda,\theta}),\qquad t=0,\ldots,T-1.
    \end{cases}
\end{align}
The perturbed objective is
\begin{align}
    J^\lambda(\theta) = \E\left[\sum_{t=0}^{T-1} r_t(X_t^{\lambda,\theta}, M_t^{\lambda,\theta}, \alpha_t^{\lambda,\theta}) + g(X_T^{\lambda,\theta}, M_T^{\lambda,\theta}) \right],
\end{align}
where the expectation is taken over the perturbations, the state process and the randomized actions.

\subsection{Density and score of the randomized perturbation}

For each $t$, let $U_t$ be an $\R^{N-1}$-valued random vector with density $\rho$ with respect to the Lebesgue measure on $\R^{N-1}$. Define $q_t = \varphi(U_t)\in\Delta_N^\circ$, with the map $\varphi : \R^{N-1}\to \Delta_N^\circ$ given by
\begin{align}
    \varphi(u) = \left( 
        \frac{\exp(u_1)}{1 + \sum_{j=1}^{N-1}\exp(u_j)},\ldots, \frac{\exp(u_{N-1})}{1 + \sum_{j=1}^{N-1}\exp(u_j)}, \frac{1}{1 + \sum_{j=1}^{N-1}\exp(u_j)}
    \right).
\end{align}
Its inverse is
\begin{align}
    \varphi^{-1}(p) = \left(\log\frac{p_1}{p_N},\ldots,\log\frac{p_{N-1}}{p_N} \right),\qquad p\in\Delta_N^\circ.
\end{align}

For fixed $(t,\lambda,\theta)$, define the affine map $\phi_t^{\lambda, \theta}:\Delta_N^\circ\to\Delta_N^\circ$ such that $\phi_t^{\lambda,\theta}(p)\coloneqq (1-\lambda)\mu_t^\theta + \lambda p$. Then
\begin{align}
    M_t^{\lambda,\theta}(\omega) = \phi_t^{\lambda,\theta}(\varphi(U_t(\omega))),\qquad
    \Lc(M_t^{\lambda,\theta}) = (\phi_t^{\lambda,\theta} \circ \varphi)\sharp \Lc(U_t).
\end{align}

\begin{theorem}[Perturbation estimate]
    Let $\lambda\in(0,1)$ and let $q_t\in\Delta_N$. For the perturbed population law $M_t^{\lambda,\theta}(\omega) = (1-\lambda)\mu_t^\theta + \lambda q_t(\omega)$, we have
    \begin{align}
        d_\mathrm{TV}(M_t^{\lambda,\theta}(\omega), \mu_t^\theta) \leq \lambda, \qquad \text{for every }\omega.
    \end{align}
    Moreover, if $\Xc$ is endowed with a metric $d_\Xc$ of diameter $\mathrm{diam}(\Xc)\coloneqq \max_{i,j\in\Xc} d_\Xc(i,j)$, then
    \begin{align}
        W_1(M_t^{\lambda,\theta}(\omega), \mu_t^\theta)\leq \mathrm{diam}(\Xc)\lambda,\qquad \text{for every }\omega.
    \end{align}
    Consequently, for any of the distances above, there exists a function $h(\lambda)$ such that
    \begin{align}
        \E\left[d(M_t^{\lambda,\theta}, \mu_t^\theta) \right]\leq h(\lambda),\qquad h(\lambda)\xrightarrow[\lambda\downarrow0]{}0.
    \end{align}
\end{theorem}

\begin{proof}
    For every $\omega$, $M_t^{\lambda,\theta}(\omega)-\mu_t^\theta = \lambda(q_t(\omega)-\mu_t^\theta)$, hence $\| M_t^{\lambda,\theta}(\omega)-\mu_t^\theta\|_1 = \lambda\|q_t(\omega)-\mu_t^\theta\|_1\leq 2\lambda$ because both $q_t$ and $\mu_t^\theta$ belong to $\Delta_N$. The total variance estimate follows from $d_\mathrm{TV}(p,p')=\frac{1}{2}\|p-p'\|_1$ on a finite space \cite{levin2026markov}.

    For the Wasserstein estimate, let $\gamma$ be any coupling of $M_t^{\lambda,\theta}$ and $\mu_t^\theta$. Since $d_\Xc(i,j)\leq\mathrm{diam}(\Xc)$, we have $\sum_{i,j\in\Xc} d_\Xc(i,j)=\gamma(i,j)\leq \mathrm{diam}(\Xc)\sum_{i\neq j}\gamma(i,j)$. Taking the infimum over all couplings:
    \begin{align}
        W_1(M_t^{\lambda,\theta}(\omega),\mu_t^\theta)\leq\mathrm{diam}(\Xc)d_\mathrm{TV}(M_t^{\lambda,\theta},\mu_t^\theta).
    \end{align}
    Thus $W_1(M_t^{\lambda,\theta}, \mu_t^\theta) \leq \mathrm{diam}(\Xc)\lambda$. The expectation estimate follows immediately since the previous bounds hold almost surely.
\end{proof}

We can now compute and get an explicit score function of the density of $M_t^{\lambda,\theta}$ on the simplex.

\begin{proposition}[Density of the logistic perturbation]
    Assume that $U_t$ has density $\rho$ on $\R^{N-1}$. Then $q_t = \varphi(U_t)$ admits the density
    \begin{align}
        f_q(p) = \rho\left(\log\frac{p_1}{p_N},\ldots,\log\frac{p_{N-1}}{p_N} \right) \frac{1}{\prod_{i=1}^N p_i}, \qquad p\in\Delta_N^\circ.
    \end{align}
\end{proposition}

The proof is a change-of-variables calculation on the simplex coordinates and is given in Appendix~\ref{appendix:discrete-density-proofs}.

\begin{proposition}[Density of the perturbed population law]
    Let $\Lambda_t^{\lambda,\theta}\coloneqq \{(1-\lambda)\mu_t^\theta+\lambda m, m\in\Delta_N^\circ\}$. For $m\in\Lambda_t^{\lambda,\theta}$, define $p_i^{\lambda,\theta}(m)\coloneqq (m_i - (1-\lambda)\mu_t^\theta(i))/\lambda$ for $i=1,\ldots, N-1$ and $p_N^{\lambda,\theta}(m)=(\phi_t^{\lambda,\theta})^{-1}(m)$. Then $M_t^{\lambda,\theta}$ admits the density
    \begin{align}
        f_t^{\lambda,\theta}(m)
        = \frac{1}{\lambda^{N-1}}
        \rho\left(\log\frac{p_1^{\lambda,\theta}(m)}{p_N^{\lambda,\theta}(m)},\ldots,\log\frac{p_{N-1}^{\lambda,\theta}(m)}{p_N^{\lambda,\theta}(m)} \right)
        \frac{1}{\prod_{i=1}^N p_i^{\lambda,\theta}(m)}
        \mathbf 1_{\Lambda_t^{\lambda,\theta}}(m),
    \end{align}
    and $\Lc(M_t^{\lambda,\theta})(\d m) = f_t^{\lambda,\theta}(m)\d m$, where $\d m$ denotes the Lebesgue measure on the simplex in the coordinates $(m_1,\ldots,m_{N-1})$.
\end{proposition}

The proof follows by applying the affine change of variables after the logistic transformation; see Appendix~\ref{appendix:discrete-density-proofs}.

\begin{corollary}[Score of the perturbation density]
    \label{corollary:discrete-score-density}
    Assume that $\rho$ is strictly positive and continuously differentiable on $\R^{N-1}$. Let $m\in\Lambda_t^{\lambda,\theta}$. Write $p_i\coloneqq p_i^{\lambda,\theta}(m)$, and set $z_i\coloneqq \log\frac{p_i}{p_N}$ and $a_i(z)\coloneqq\partial_{z_i}\log\rho(z)$ for $i=1,\ldots,N-1$. Then for fixed $m$ in the interior of the support,
    \begin{align}
        \nabla_\theta \log f_t^{\lambda,\theta}(m) = -\frac{1-\lambda}{\lambda}\sum_{k=1}^{N-1} H_k(p) \nabla_\theta \mu_t^\theta(k),
    \end{align}
    where $H_k(p) = \frac{a_k(z)}{p_k} + \frac{\sum_{i=1}^{N-1} a_i(z)}{p_N} - \frac{1}{p_k} + \frac{1}{p_N}$ for $k=1,\ldots,N-1$.
\end{corollary}

The proof differentiates the previous density at fixed $m$ and collects the coefficients of the independent coordinates of $\nabla_\theta\mu_t^\theta$; see Appendix~\ref{appendix:discrete-density-proofs}.

\subsection{Convergence of the perturbed objective}

We now show that the perturbed control problem is an approximation of the original one when $\lambda$ is small. The law of the perturbed state process, $\nu_t^{\lambda,\theta}\coloneqq \Lc(X_t^{\lambda,\theta})$, needs to be shown to remain close to $\mu_t^\theta$.

For $\theta\in\Theta$, define the averaged transition kernel, a probability measure on $\Xc$ and the conditional law of the next state $X_{t+1}$ given $X_t = i$, after averaging over the randomized action $\alpha_t\sim\pi_t^\theta(\cdot\mid i,m)$.
\begin{align}
    K_t^\theta(j\mid i,m) \coloneqq \int_\Ac P_t(j\mid i,m,a) \pi_t^\theta(\d a\mid i,m),\qquad i,j\in\Xc,\ m\in\Delta_N,
\end{align}
and the averaged running reward
\begin{align}
    R_t^\theta(i,m)\coloneqq \int_\Ac r_t(i,m,a)\pi_t^\theta(\d a\mid i,m),\qquad i\in\Xc,\ m\in\Delta_N.
\end{align}
The unperturbed law and the perturbed dynamics each then satisfy
\begin{align}
    \mu_{t+1}^\theta(j) = \sum_{i=1}^N \mu_t^\theta(i) K_t^\theta(j\mid i,\mu_t^\theta),\qquad
    \nu_{t+1}^{\lambda,\theta}(j) = \E\left[K_t^\theta(j\mid X_t^{\lambda,\theta}, M_t^{\lambda,\theta}) \right],\qquad
    j\in\Xc.
\end{align}

Indeed, by the tower property,
\begin{align}
    \nu_{t+1}^{\lambda,\theta}(j)
    &= \P(X_{t+1}^{\lambda,\theta}=j) = \E\left[\mathbf 1_{\{X_{t+1}^{\lambda,\theta}=j \}} \right] \\
    &= \E\left[\E\left[\mathbf 1_{\{X_{t+1}^{\lambda,\theta}=j \}} \mid X_t^{\lambda,\theta}, M_t^{\lambda,\theta}, \alpha_t^{\lambda,\theta}\right] \right] \\
    &= \E\left[P_t(j\mid X_t^{\lambda,\theta}, M_t^{\lambda,\theta}\alpha_t^{\lambda,\theta}) \right].
\end{align}
Conditioning now on $(X_t^{\lambda,\theta}, M_t^{\lambda,\theta})$, and using $\alpha_t^{\lambda,\theta} \sim \pi_t^\theta(\cdot\mid X_t^{\lambda,\theta}, M_t^{\lambda,\theta})$, we obtain
\begin{align}
    \nu_{t+1}^{\lambda,\theta}
    &= \E\left[\int_A P_t(j\mid X_t^{\lambda,\theta}, M_t^{\lambda,\theta}, a) \pi_t^\theta(\d a\mid X_t^{\lambda,\theta},M_t^{\lambda,\theta}) \right] \\
    &= \E\left[K_t^\theta(j\mid X_t^{\lambda,\theta}, M_t^{\lambda,\theta}) \right].
\end{align}

\begin{assumption}[Uniform regularity]
    \label{ass:discrete-uniform-regularity}
    There exists constants $L_K,L_R,L_g\geq0$ and $\bar R,\bar g<\infty$ such that, for every $\theta\in\Theta,\ t=0,\ldots,T-1,\ i\in\Xc$, and $m,m'\in\Delta_N$,
    \begin{align}
        d_\mathrm{TV}\left(K_t^\theta(\cdot\mid i,m),\ K_t^\theta(\cdot\mid i,m') \right)
        &\leq L_K d_\mathrm{TV}(m,m'), \\
        |R_t^\theta(i,m) - R_t^\theta(i,m')|
        &\leq L_R d_\mathrm{TV}(m,m'), \\
        |g(i,m) - g(i,m')|
        &\leq L_g d_\mathrm{TV}(m,m'), \\
        |R_t^\theta(i,m)|\leq \bar R,&\quad |g(i,m)|\leq\bar g.
    \end{align}
\end{assumption}

\begin{lemma}[Stability of the state marginal]
    \label{lemma:stability-state-marginal}
    Under Assumption~\ref{ass:discrete-uniform-regularity}, for every $\theta\in\Theta$ and every $t=0,\ldots,T$,
    \begin{align}
        d_\mathrm{TV}(\nu_t^{\lambda,\theta},\mu_t^\theta)\leq L_K\lambda t.
    \end{align}
\end{lemma}

The proof is a finite-horizon induction based on the dual representation of total variation; see Appendix~\ref{appendix:objective-convergence-proofs}.

\begin{theorem}[Convergence of the perturbed objective]
    \label{theorem:discrete-convergence-objective}
    Under Assumption~\ref{ass:discrete-uniform-regularity}, for every $\theta\in\Theta$,
    \begin{align}
        |J^\lambda(\theta) - J(\theta)|\leq C_T\lambda.
    \end{align}
    If the constants in the assumption are uniform in $\theta$, the the convergence is uniform.
\end{theorem}

The proof decomposes running and terminal reward errors into marginal and population-argument errors, then uses Lemma~\ref{lemma:stability-state-marginal}; see Appendix~\ref{appendix:objective-convergence-proofs}.

\begin{corollary}[Consistency of optimizers]
    Assume that $\Theta$ is compact, that $J$ is continuous on $\theta$, and that $\sup_{\theta\in\Theta} |J^\lambda(\theta) - J(\theta)|\xrightarrow[\lambda\downarrow0]{}0$. Let $\theta_\lambda^\star\in\argmax_{\theta\in\Theta} J^\lambda(\theta)$ and suppose that, along a sequence $\lambda_n\downarrow0$, $\theta_{\lambda_n}^\star \to \bar \theta$. Then
    \begin{align}
        \bar \theta \in \argmax_{\theta\in\Theta}J(\theta).
    \end{align}
    In particular, if $J$ admits a unique maximizer $\theta^\star$, then $\theta_\lambda^\star \to \theta^\star$ as $\lambda\downarrow0$.
\end{corollary}

This is the standard argmax-continuity argument under uniform convergence; see Appendix~\ref{appendix:optimizer-consistency-proof}.

\subsection{Gradient-level convergence}

Theorem~\ref{theorem:discrete-convergence-objective} proves the convergence of the perturbed analysis, but a stronger gradient-level version is needed for the statistical analysis of the policy-gradient estimator. This requires additional smoothness assumptions. We use the notation $D_t^\theta[h]\coloneqq \partial_h\mu_t^\theta$ and $E_t^{\lambda,\theta}[h]\coloneqq\partial_h \nu_t^{\lambda,\theta}$ for the directional derivatives in the direction $h\in\R^{d_\theta}$, where $d_\theta$ is the dimension of the parameter space. We assume $|h|=1$.

\begin{assumption}[Smoothness of the averaged dynamics]
    \label{ass:discrete-gradient-regularity}
    The maps $(\theta,m)\mapsto K_t^\theta(\cdot\mid i,m)$ and $(\theta,m)\mapsto R_t^\theta(i,m)$ are continuously differentiable with respect to $\theta$ and $m$, uniformly in $t$ and $i$. The terminal reward $m\mapsto g(i,m)$ is continuously differentiable.

    Moreover, there exists a constant $C_1 > 0$ such that, for every unit vector $h$, every vector $v\in\R^N$ satisfying $\sum_i v_i=0$, and every $m,m'\in\Delta_N$,
    \begin{align}
        d_\mathrm{TV}\left(\partial_h K_t^\theta(\cdot\mid i,m), \partial_h K_t^\theta(\cdot\mid i,m') \right) &\leq C_1 d_\mathrm{TV}(m,m'), \\
        d_\mathrm{TV}\left(D_m K_t^\theta(\cdot\mid i,m)[v], D_m K_t^\theta(\cdot\mid i,m')[v] \right) &\leq C_1\|v\|_1 d_\mathrm{TV}(m,m'), \\
        d_\mathrm{TV}\left(D_m K_t^\theta(\cdot\mid i,m)[v], D_m K_t^\theta(\cdot\mid i,m)[v'] \right) &\leq C_1 \|v-v'\|_1.
    \end{align}
    Similarly,
    \begin{align}
        \left|\partial_h R_t^\theta(i,m) - \partial_h R_t^\theta(i,m') \right| &\leq C_1d_\mathrm{TV}(m,m'), \\
        \left|D_m R_t^\theta(i,m)[v] - D_mR_t^\theta(i,m')[v] \right| &\leq C_1\|v\|_1 d_\mathrm{TV}(m,m'), \\
        \left|D_m R_t^\theta(i,m)[v] - D_mR_t^\theta(i,m)[v'] \right| &\leq C_1\|v-v'\|_1.
    \end{align}
    The same type of bounds hold for $g$, namely
    \begin{align}
        \left| D_mg(i,m)[v] - D_mg(i,m')[v]\right| &\leq C_1\|v\|_1 d_\mathrm{TV}(m,m'),\\
        \left| D_mg(i,m)[v] - D_mg(i,m)[v']\right| &\leq C_1\|v -v'\|_1.
    \end{align}
    Finally, all first-order derivations appearing above are uniformly bounded.
\end{assumption}

\begin{lemma}[Stability of the differentiated state marginal]
    Under Assumption~\ref{ass:discrete-gradient-regularity}, there exists a constant $C_T^D > 0$ such that, for every $\theta\in\Theta$, every $\lambda\in(0,1)$, every unit direction $h$, and every $t=0,\ldots,T$,
    \begin{align}
        \left\|E_t^{\lambda,\theta}[h] - D_t^\theta[h] \right\|_1 \leq C_T^D\lambda.
    \end{align}
\end{lemma}

The proof differentiates the two marginal recursions, compares the resulting triangular systems, and applies a finite-horizon induction; see Appendix~\ref{appendix:gradient-convergence-proofs}.

\begin{theorem}[Gradient-level convergence]
    \label{theorem:discrete-gradient-convergence}
    Under Assumption~\ref{ass:discrete-gradient-regularity}, there exists a constant $C_T^\nabla > 0$ such that, for every $\theta\in\Theta$ and every $\lambda\in(0,1)$,
    \begin{align}
        \|\nabla_\theta J^\lambda(\theta) - \nabla_\theta J(\theta)\| \leq C_T^\nabla \lambda.
    \end{align}
    If the constants in the assumption are uniform in $\theta$, then the convergence is uniform in $\theta$.
\end{theorem}

The proof compares directional derivatives of $J^\lambda$ and $J$ term by term and uses the differentiated marginal stability estimate; see Appendix~\ref{appendix:gradient-convergence-proofs}.

\subsection{Policy-gradient formula for the perturbed problem}

We now derive a likelihood-ratio representation for the gradient of the perturbed objective $J^\lambda$. The perturbation introduced transfers the dependence on $\theta$ through the population law to the density $f_t^{\lambda,\theta}$ of the random variable $M_t^{\lambda,\theta}$. Differentiating $J^\lambda$ therefore produces a standard policy-score term and an additional score term coming from the density of the perturbation.

Define the perturbed return
\begin{align}
    \Rc_\theta^\lambda \coloneqq \sum_{t=0}^{T-1} r_t(X_t^{\lambda,\theta}, M_t^{\lambda,\theta}, \alpha_t^{\lambda,\theta}) + g(X_T^{\lambda,\theta}, M_T^{\lambda,\theta}).
\end{align}

The complete trajectory of the perturbed system is denoted by
\begin{align}
    \Gamma^\lambda \coloneqq \left(M_0^{\lambda,\theta}, X_0^{\lambda,\theta}, \alpha_0^{\lambda,\theta}, \ldots, M_{T-1}^{\lambda,\theta}, X_{T-1}^{\lambda,\theta}, \alpha_{T-1}^{\lambda,\theta}, M_T^{\lambda,\theta}, X_T^{\lambda,\theta} \right).
\end{align}
The law of this trajectory has density, with respect to the product reference measure, proportional to
\begin{align}
    \mu_0(x_0) \prod_{t=0}^{T-1}\left[f_t^{\lambda,\theta}(m_t) p_t^\theta(a_t\mid x_t, m_t) P_t(x_{t+1}\mid x_t,a_t,m_t) \right] f_T^{\lambda,\theta}(m_T).
\end{align}
Here, the transition probabilities $P_t$ do not depend directly on $\theta$. Their dependence on $\theta$ has been moved to the random population argument $M_t^{\lambda,\theta}$ through its density $f_t^{\lambda,\theta}$.

\begin{assumption}[Differentiability and integrability]
    For every $t=0,\ldots,T-1$, the randomized policy admits a density with respect to a reference measure $\nu_\Ac$ on $\Ac$:
    \begin{align}
        \pi_t^\theta(\d a\mid i,m) = p_t^\theta(a\mid i,m)\nu_\Ac(\d a),\qquad i\in\Xc,\ m\in\Delta_N.
    \end{align}
    Assume additionally that $p_t^\theta(a\mid i,m) > 0$ and that $p_t^\theta(a\mid i,m)$ is continuously differentiable with respect to $\theta$. With a slight abuse of notation, we write $\nabla_\theta \log\pi_t^\theta(a\mid i,m)\coloneqq \nabla_\theta\log p_t^\theta(a\mid i,m)$.

    Moreover, for every $t=0,\ldots,T$, the perturbation density $f_t^{\lambda,\theta}$ is strictly positive on the relative interior of its support $\Lambda_t^{\lambda,\theta}$, and is continuously differentiable with respect to $\theta$ for fixed $m$ in this support.

    Finally, we assume that differentiation under the integral sign is justified. Equivalently, the trajectory score
    \begin{align}
        S_\theta^\lambda = \sum_{t=0}^{T-1} \nabla_\theta \log p_t^\theta(\alpha_t^{\lambda,\theta}\mid X_t^{\lambda,\theta}, M_t^{\lambda,\theta}) + \sum_{t=0}^T \nabla_\theta \log f_t^{\lambda,\theta}(M_t^{\lambda,\theta})
    \end{align}
    is integrable against the perturbed return, namely $\E[\|\Rc_\theta^\lambda S_\theta^\lambda\|] < \infty$.
\end{assumption}

\begin{theorem}[Policy-gradient formula for the perturbed objective]
    \label{theorem:discrete-policy-gradient-formula}
    Under the previous assumptions,
    \begin{align}
        \nabla_\theta J^\lambda(\theta) = \E[\Rc_\theta^\lambda S_\theta^\lambda],
    \end{align}
    where, using Corollary~\ref{corollary:discrete-score-density}, the score of the perturbed trajectory is
    \begin{align}
        S_\theta^\lambda
        &\coloneqq \sum_{t=0}^{T-1} \nabla_\theta \log p_t^\theta(\alpha_t^{\lambda,\theta}\mid X_t^{\lambda,\theta}, M_t^{\lambda,\theta}) + \sum_{t=0}^T \nabla_\theta \log f_t^{\lambda,\theta}(M_t^{\lambda,\theta})\\
        &= \sum_{t=0}^{T-1} \nabla_\theta \log p_t^\theta(\alpha_t^{\lambda,\theta}\mid X_t^{\lambda,\theta}, M_t^{\lambda,\theta}) - \frac{1-\lambda}{\lambda} \sum_{t=0}^T \sum_{k=1}^{N-1} H_k(q_t) \nabla_\theta \mu_t^\theta(k).
    \end{align}
\end{theorem}

\begin{proof}
    Let $L_\theta^\lambda$ be the density of the complete perturbed trajectory $\Gamma^\lambda$. The only direct dependence of $L_\theta^\lambda$ on $\theta$ comes from the policy densities $p_t^\theta(a_t\mid x_t,m_t)$ and from the perturbation densities $f_t^{\lambda,\theta}(m_t)$; the transition kernels have no direct $\theta$-dependence once the sampled population arguments are fixed. Therefore
    \begin{align}
        \nabla_\theta\log L_\theta^\lambda(\gamma)
        =
        \sum_{t=0}^{T-1}\nabla_\theta\log p_t^\theta(a_t\mid x_t,m_t)
        +
        \sum_{t=0}^T\nabla_\theta\log f_t^{\lambda,\theta}(m_t)
        =
        S_\theta^\lambda(\gamma).
    \end{align}
    Differentiating under the integral sign and using $\nabla_\theta L_\theta^\lambda=L_\theta^\lambda\nabla_\theta\log L_\theta^\lambda$ yields
    \begin{align}
        \nabla_\theta J^\lambda(\theta)
        &=
        \nabla_\theta\int \Rc^\lambda(\gamma)L_\theta^\lambda(\gamma)\,\d\gamma
        =
        \int \Rc^\lambda(\gamma)\nabla_\theta L_\theta^\lambda(\gamma)\,\d\gamma\\
        &=
        \int \Rc^\lambda(\gamma)S_\theta^\lambda(\gamma)L_\theta^\lambda(\gamma)\,\d\gamma
        =
        \E[\Rc_\theta^\lambda S_\theta^\lambda].
    \end{align}
\end{proof}

\begin{corollary}[Baseline version]
    Let $b\in\R$ be any deterministic baseline. Then
    \begin{align}
        \nabla_\theta J^\lambda(\theta) = \E[(\Rc_\theta^\lambda - b)S_\theta^\lambda].
    \end{align}
    More generally, one may use any baseline $b$ such that $\E[bS_\theta^\lambda] = 0$.
\end{corollary}

\begin{proof}
    Since $L_\theta^\lambda$ is a probability density, differentiating $\int L_\theta^\lambda(\gamma)\,\d\gamma=1$ gives $\E[S_\theta^\lambda]=0$. Hence, for any deterministic $b$,
    \begin{align}
        \E[\Rc^\lambda_\theta S^\lambda_\theta]
        =
        \E[(\Rc^\lambda_\theta-b)S^\lambda_\theta]+b\E[S^\lambda_\theta]
        =
        \E[(\Rc^\lambda_\theta-b)S^\lambda_\theta].
    \end{align}
\end{proof}


\subsection{Estimation of the population-flow sensitivity}

The policy-gradient formula from Theorem~\ref{theorem:discrete-policy-gradient-formula} still contains the population-flow sensitivity $D_t^\theta(k)\coloneqq \nabla_\theta \mu_t^\theta(k)$, for $t=0,\ldots,T$ and $k=1,\ldots,N-1$. This quantity is not directly observed from a single trajectory. Moreover, the usual likelihood-ratio identity does not provide an exact model-free estimator in the mean-field setting, because the transition kernel and the policy depend on the whole population law $\mu_t^\theta$, which itself depends on $\theta$.

We therefore construct an auxiliary estimator of $D_t^\theta$ using the same perturbation idea as above. Since $\sum_{i=1}^N \mu_t^\theta(i) = 1$, we have $\sum_{i=1}^N D_t^\theta(i) = 0$ and it is enough to estimate $D_t^\theta(1),\ldots, D_t^\theta(N-1)$ and then recover $D_t^\theta(N)$.

Let $\eta\in(0,1)$ be an auxiliary perturbation parameter, possibly different from the perturbation $\lambda$ used in the policy-gradient estimator. For $s=0,\ldots,T$, let $\bar M_s^{\eta,\theta} \coloneqq (1-\eta)\mu_s^\theta + \eta \bar q_s$, where $\bar q_s\in\Delta_N^\circ$ is generated by the same logistic parametrization as before. The auxiliary perturbed process is
\begin{align}
    \begin{cases}
        \bar{X}_0^{\eta, \theta} \sim \mu_0, \\
        \bar{\alpha}_s^{\eta, \theta} \sim \pi_s^\theta(\cdot \mid \bar{X}_s^{\eta, \theta}, \bar{M}_s^{\eta, \theta}), \\
        \bar{X}_{s+1}^{\eta, \theta} \sim P_s(\cdot \mid \bar{X}_s^{\eta, \theta}, \bar{M}_s^{\eta, \theta}, \bar\alpha_s^{\eta, \theta}), \qquad s=0,\ldots, T-1.
    \end{cases}
\end{align}
The corresponding policy-score term is $\Psi_s^{\eta, \theta} \coloneqq \nabla_\theta \log p_s^\theta (\bar\alpha_s^{\eta, \theta} \mid \bar X_s^{\eta, \theta}, \bar M_s^{\eta, \theta})$. The perturbation-density score at time $s$ is
\begin{align}
    -\frac{1-\eta}{\eta} \sum_{\ell=1}^{N-1}H_\ell(\bar q_s)D_s^\theta(\ell).
\end{align}

For fixed $t\in\{1,\ldots,T\}$ and $k\in\{1,\ldots,N-1\}$, consider the auxiliary objective $\mu_{t,k}^\eta(\theta) \coloneqq \E\left[\mathbf 1_{\{\bar X_t^{\eta,\theta}=k \}} \right]$. Applying the perturbed likelihood ratio formula to the terminal reward $\mathbf 1_{\{\bar X_t^{\eta,\theta}=k \}}$ and using causality, gives the formal identity
\begin{align}
    \nabla_\theta \mu_{t,k}^\eta(\theta) = \E \left[ \mathbf{1}_{\{\bar X_t^{\eta, \theta}=k\}} \sum_{s=0}^{t-1}\Xi_s^{\eta, \theta} \right].
\end{align}
The terms with $s\geq t$ vanish because the event $\{\bar X_t^{\eta, \theta}=k\}$ does not depend on future actions or future perturbations.

Since $\bar M_s^{\eta, \theta}$ converges to $\mu_s^\theta$ as $\eta\downarrow0$, we expect $\mu_{t,k}^\eta(\theta) \xrightarrow[\eta\downarrow0]{} \mu_t^\theta(k)$. Under a gradient-level convergence assumption, this yields
\begin{align}
    D_t^\theta(k) = \lim_{\eta\downarrow0}\E\left[ \mathbf{1}_{\{\bar X_t^{\eta,\theta}=k\}} \sum_{s=0}^{t-1} \left(\Psi_s^{\eta, \theta} - \frac{1-\eta}{\eta}\sum_{\ell=1}^{N-1}H_\ell(\bar q_s)D_s^\theta(\ell) \right) \right].
\end{align}
This relation is triangular in time: the sensitivity $D_t^\theta$ depends only on $D_0^\theta, \ldots, D_{t-1}^\theta$. Since the initial law $\mu_0$ is independent of $\theta$, we have $D_0^\theta(k) = 0$ for $k=1,\ldots,N$. Therefore, the sensitivities can be estimated recursively forward in time.

Let $n\geq 1$ be the number of auxiliary trajectories used to estimate the population-flow sensitivity. For each $r=1,\ldots,n$, simulate an independent auxiliary perturbed trajectory $\left(\bar q_s^{(r)}, \bar M_s^{\eta,\theta,(r)}, \bar X_s^{\eta,\theta,(r)}, \bar \alpha_s^{\eta,\theta,(r)} \right)_{s=0}^T$. Define $\psi_s^{\eta,\theta,(r)} \coloneqq \nabla_\theta \log p_s^\theta(\bar \alpha_s^{\eta,\theta,(r)}\mid \bar X_s^{\eta,\theta,(r)}, \bar M_s^{\eta,\theta,(r)})$. We initialize $\widehat D_0^{\eta,n,\theta}(k)=0$ for $k=1,\ldots,N$. Then, for $t=1,\ldots,T$ and $k=1,\ldots,N-1$, define recursively
\begin{align}
    \widehat D_t^{\eta,n,\theta} (k) \coloneqq \frac{1}{n} \sum_{r=1}^n \mathbf 1_{\{\bar X_t^{\eta,\theta,(r)}=k\}} \left[\sum_{s=0}^{t-1}\psi_s^{\eta,\theta,(r)} - \frac{1-\eta}{\eta}\sum_{s=0}^{t-1} \sum_{\ell=1}^{N-1} H_\ell(\bar q_s^{(r)}) \widehat D_s^{\eta,n,\theta}(\ell) \right].
\end{align}
Finally, set $\widehat D_t^{\eta,n,\theta}(N) \coloneqq - \sum_{k=1}^{N-1} \widehat D_t^{\eta,n,\theta}(k)$.

\begin{remark}[Empirical population law]
    In a fully simulation-based implementation, in order to construct $\bar M_s^{\eta,\theta} = (1-\eta)\mu_s^\theta + \eta\bar q_s$, $\mu_s^\theta$ may be replaced by an empirical population law
    \begin{align}
        \widehat \mu_s^\theta(i) = \frac{1}{M} \sum_{m=1}^M \mathbf 1_{\{X_s^{\theta,(m)} = i\}},\qquad i=1,\ldots,N,
    \end{align}
    computed from a batch of independent unperturbed trajectories. This introduces an additional sampling error, which can be treated separately in the final mean-square error decomposition.
\end{remark}

\subsection{Monte Carlo estimator of the perturbed policy gradient}

We now combine the perturbed policy-gradient formula with the auxiliary estimator of the population-flow sensitivity to obtain a fully sample-based estimator of $\nabla_\theta J^\lambda (\theta)$. The only non-directly observable term in the policy-gradient formula is $\left(D_t^\theta(k) \right)_{0\leq t\leq T,\ 1\leq k\leq N-1}$, which we replace by the auxiliary estimator $\widehat D_t^{\eta,n,\theta}(k)$.

Let $B \geq 1$ be the number of trajectories used to estimate the perturbed policy gradient, and assume that the auxiliary batch used to compute $\widehat D_t^{\eta,n,\theta}(k)$ is independent of the batch used below. For each $b=1,\ldots,B$, simulate an independent perturbed trajectory $\left(q_t^{(b)}, M_t^{\lambda,\theta,(b)}, X_t^{\lambda,\theta,(b)},\alpha_t^{\lambda,\theta,(b)} \right)_{t=0}^T$ according to
\begin{align}
    \begin{cases}
        X_0^{\lambda,\theta,(b)} &\sim \mu_0,\\
        M_t^{\lambda,\theta,(b)} &=(1-\lambda)\mu_t^\theta + \lambda q_t^{(b)},\\
        \alpha_t^{\lambda,\theta,(b)} &\sim \pi_t^\theta(\cdot\mid X_t^{\lambda,\theta,(b)}, M_t^{\lambda,\theta,(b)}),\\
        X_{t+1}^{\lambda,\theta,(b)} &\sim P_t(\cdot\mid X_t^{\lambda,\theta,(b)}, M_t^{\lambda,\theta,(b)}, \alpha_t^{\lambda,\theta,(b)}).
    \end{cases}
\end{align}
The corresponding perturbed return is
\begin{align}
    \Rc_\theta^{\lambda,(b)} \coloneqq \sum_{t=0}^{T-1} r_t(X_t^{\lambda,\theta,(b)}, M_t^{\lambda,\theta,(b)},\alpha_t^{\lambda,\theta,(b)}) + g(X_T^{\lambda,\theta,(b)},M_T^{\lambda,\theta,(b)}).
\end{align}
The policy-score term is
\begin{align}
    \Sc_\mathrm{pol}^{\lambda,\theta,(b)} \coloneqq \sum_{t=0}^{T-1} \nabla_\theta \log p_t^\theta(\alpha_t^{\lambda,\theta,(b)} \mid X_t^{\lambda,\theta,(b)}, M_t^{\lambda,\theta,(b)}).
\end{align}
It is useful to isolate the ideal estimator, in which the exact population-flow sensitivity is available. Define
\begin{align}
    \Sc_\mathrm{pert}^{\lambda,\theta,(b)}
    &\coloneqq
    -\frac{1-\lambda}{\lambda}
    \sum_{t=0}^T\sum_{k=1}^{N-1}
    H_k(q_t^{(b)})D_t^\theta(k),\\
    \Sc^{\lambda,\theta,(b)}
    &\coloneqq
    \Sc_\mathrm{pol}^{\lambda,\theta,(b)}
    +
    \Sc_\mathrm{pert}^{\lambda,\theta,(b)}.
\end{align}
The corresponding exact-sensitivity, or oracle, Monte Carlo estimator is
\begin{align}
    G_{\lambda,B}(\theta)
    \coloneqq
    \frac1B
    \sum_{b=1}^B
    \left(\Rc_\theta^{\lambda,(b)}-b_0\right)
    \Sc^{\lambda,\theta,(b)}.
\end{align}
By Theorem~\ref{theorem:discrete-policy-gradient-formula} and the baseline identity, $G_{\lambda,B}(\theta)$ is exactly unbiased for $\nabla_\theta J^\lambda(\theta)$. This estimator is not model-free, since it assumes access to the exact sensitivities $D_t^\theta(k)$, but it is a useful oracle reference: it isolates the effect of the transport perturbation from the additional error caused by estimating the mean-field derivative.
The perturbation-score term is estimated by
\begin{align}
    \widehat \Sc_\mathrm{pert}^{\lambda,\eta,n,\theta,(b)} \coloneqq -\frac{1-\lambda}{\lambda} \sum_{t=0}^T \sum_{k=1}^{N-1} H_k(q_t^{(b)}) \widehat D_t^{\eta,n,\theta}(k).
\end{align}
Thus the estimated total score is $\widehat \Sc^{\lambda,\eta,n,\theta,(b)} \coloneqq \Sc_\mathrm{pol}^{\lambda,\theta,(b)} + \widehat \Sc_\mathrm{pert}^{\lambda, \eta,n,\theta,(b)}$. For a deterministic baseline $b_0\in\R$, define the Monte Carlo estimator
\begin{align}
    \widehat G_{\lambda,\eta,B,n}(\theta) \coloneqq \frac1B\sum_{b=1}^B\left(\Rc_\theta^{\lambda,(b)}-b_0\right)\widehat \Sc^{\lambda,\eta,n,\theta,(b)}.
\end{align}
The role of the baseline is only variance reduction: when the exact score is used it does not change the expectation, and in the plug-in estimator it only multiplies the same score approximation.

\subsection{Bias and mean-square error decomposition}

The oracle estimator $G_{\lambda,B}$ only contains the perturbation bias and the ordinary Monte Carlo variance. The practical estimator $\widehat G_{\lambda,\eta,B,n}$ adds a third source of error: the auxiliary error made when $D_t^\theta$ is replaced by $\widehat D_t^{\eta,n,\theta}$.

\begin{proposition}[Unbiasedness and error of the oracle estimator]
    \label{prop:oracle-mse-decomposition}
    Under the assumptions of Theorem~\ref{theorem:discrete-gradient-convergence}, assume in addition that the perturbed returns and scores have bounded second moments, with the perturbation-score contribution scaling as $O(\lambda^{-1})$. Then
    \begin{align}
        \E[G_{\lambda,B}(\theta)]
        &=
        \nabla_\theta J^\lambda(\theta),\\
        \E\left[
            \left\|
                G_{\lambda,B}(\theta)
                -
                \nabla_\theta J^\lambda(\theta)
            \right\|^2
        \right]
        &\leq
        C_T\frac{1+\lambda^{-2}}{B}.
    \end{align}
    Consequently, with respect to the original objective,
    \begin{align}
        \left\|
            \E[G_{\lambda,B}(\theta)]
            -
            \nabla_\theta J(\theta)
        \right\|
        &\leq C_T\lambda,\\
        \E\left[
            \left\|
                G_{\lambda,B}(\theta)
                -
                \nabla_\theta J(\theta)
            \right\|^2
        \right]
        &\leq
        C_T\left(
            \lambda^2
            +
            \frac{1+\lambda^{-2}}{B}
        \right).
    \end{align}
\end{proposition}

\begin{proof}
    The exact unbiasedness follows from Theorem~\ref{theorem:discrete-policy-gradient-formula} and the baseline identity. The second estimate is the variance of an average of $B$ independent oracle-score samples, bounded by the score moment assumption. The comparison with $\nabla_\theta J(\theta)$ then follows by adding and subtracting $\nabla_\theta J^\lambda(\theta)$ and using the gradient convergence estimate of Theorem~\ref{theorem:discrete-gradient-convergence}.
\end{proof}

\begin{proposition}[Auxiliary sensitivity error]
    \label{prop:auxiliary-sensitivity-error}
    Assume the gradient-level convergence estimate of Theorem~\ref{theorem:discrete-gradient-convergence} and suppose that the auxiliary likelihood-ratio scores have uniformly bounded second moments, with the perturbation-density score scaling as $O(\eta^{-1})$. Then, for every $t=0,\ldots,T$ and $k=1,\ldots,N-1$, there exists a constant $C_T>0$ such that
    \begin{align}
        \E\left[
            \left\|
                \widehat D_t^{\eta,n,\theta}(k)-D_t^\theta(k)
            \right\|^2
        \right]
        \leq
        C_T\left(\eta^2+\frac1{n\eta^2}\right).
    \end{align}
\end{proposition}

\begin{proof}
    The auxiliary identity targets the derivative of the $\eta$-perturbed marginal, so the deterministic bias is $O(\eta)$ by Theorem~\ref{theorem:discrete-gradient-convergence}. For fixed previous sensitivities, the auxiliary estimator is an average of $n$ independent likelihood-ratio variables whose score has second moment of order $1+\eta^{-2}$, hence variance $O((n\eta^2)^{-1})$. The recursive use of $\widehat D_s^{\eta,n,\theta}$ for $s<t$ propagates these errors only through a finite triangular system. An induction over $t$ gives the stated bound.
\end{proof}

To state the full decomposition, define the main-batch coefficients
\begin{align}
    A_{t,k}^{\lambda,B}(\theta)
    \coloneqq
    \frac1B\sum_{b=1}^B
    \left(\Rc_\theta^{\lambda,(b)}-b_0\right)
    H_k(q_t^{(b)}),
    \qquad
    a_{t,k}^{\lambda}(\theta)
    \coloneqq
    \E[A_{t,k}^{\lambda,1}(\theta)].
\end{align}
Then the plug-in estimator differs from the exact-sensitivity estimator by
\begin{align}
    \widehat G_{\lambda,\eta,B,n}(\theta)-G_{\lambda,B}(\theta)
    =
    -\frac{1-\lambda}{\lambda}
    \sum_{t=0}^T\sum_{k=1}^{N-1}
    A_{t,k}^{\lambda,B}(\theta)
    \left(
        \widehat D_t^{\eta,n,\theta}(k)-D_t^\theta(k)
    \right).
\end{align}
This identity makes explicit where the factor $1/\lambda$ enters the error.

\begin{theorem}[Bias and mean-square error]
    \label{theorem:discrete-mse-decomposition}
    Assume the hypotheses of Proposition~\ref{prop:auxiliary-sensitivity-error}. Assume also that the perturbed returns, policy scores and perturbation scores have bounded second moments uniformly in $\theta$, with the perturbation-score contribution scaling as $O(\lambda^{-1})$. Then
    \begin{align}
        \left\|
            \E[\widehat G_{\lambda,\eta,B,n}(\theta)]
            -
            \nabla_\theta J(\theta)
        \right\|
        &\leq
        C_T\left(
            \lambda
            +
            \frac1{\lambda}
            \sqrt{\eta^2+\frac1{n\eta^2}}
        \right),\\
        \E\left[
            \left\|
                \widehat G_{\lambda,\eta,B,n}(\theta)
                -
                \nabla_\theta J(\theta)
            \right\|^2
        \right]
        &\leq
        C_T\left(
            \lambda^2
            +
            \frac{1+\lambda^{-2}}{B}
            +
            \frac1{\lambda^2}
            \left(\eta^2+\frac1{n\eta^2}\right)
        \right).
    \end{align}
\end{theorem}

\begin{proof}
    Let $\bar G_{\lambda,\eta,n}(\theta)\coloneqq\E[\widehat G_{\lambda,\eta,B,n}(\theta)\mid \widehat D^{\eta,n,\theta}]$, where the expectation is only over the main batch. By independence of the main and auxiliary batches,
    \begin{align}
        \bar G_{\lambda,\eta,n}(\theta)
        =
        \nabla_\theta J^\lambda(\theta)
        -
        \frac{1-\lambda}{\lambda}
        \sum_{t=0}^T\sum_{k=1}^{N-1}
        a_{t,k}^{\lambda}(\theta)
        \left(
            \widehat D_t^{\eta,n,\theta}(k)-D_t^\theta(k)
        \right).
    \end{align}
    Combining this identity with Theorem~\ref{theorem:discrete-gradient-convergence} and Proposition~\ref{prop:auxiliary-sensitivity-error} gives the bias bound.

    For the mean-square error, use the conditional variance decomposition
    \begin{align}
        \E\left[
            \left\|
                \widehat G_{\lambda,\eta,B,n}(\theta)-\nabla_\theta J(\theta)
            \right\|^2
        \right]
        =
        \E\left[
            \left\|
                \widehat G_{\lambda,\eta,B,n}(\theta)-\bar G_{\lambda,\eta,n}(\theta)
            \right\|^2
        \right]
        +
        \E\left[
            \left\|
                \bar G_{\lambda,\eta,n}(\theta)-\nabla_\theta J(\theta)
            \right\|^2
        \right].
    \end{align}
    The first term is the variance of an average of $B$ independent main-batch samples, and the score moment assumption bounds it by $C_T(1+\lambda^{-2})/B$. The second term is bounded by the square of the perturbation bias plus the auxiliary plug-in error, which gives $C_T\lambda^2+C_T\lambda^{-2}(\eta^2+(n\eta^2)^{-1})$.
\end{proof}

The bound displays the usual tuning tradeoff. Decreasing $\lambda$ reduces the deterministic perturbation bias, but increases both the variance of the perturbation score and the amplification of the auxiliary sensitivity error. The auxiliary radius $\eta$ has the same structure: small $\eta$ reduces the bias of the sensitivity target, while the likelihood-ratio variance grows like $\eta^{-2}$.

\section{Continuous-state space}

\subsection{Gaussian-mixture representation and perturbation of the population law}
\label{subsec:continuous-gmm-perturbation}

We now consider the continuous-state setting $\Xc=\R^d$, and denote by $\Pc_1(\R^d)$ the set of probability measures with finite first moment. For a policy parameter $\theta\in\Theta$, let $\mu_t^\theta\coloneqq \Lc(X_t^\theta)\in\Pc_1(\R^d)$, for $t\in\Tc$, be the population law of the unperturbed controlled process.

The main difference with the finite-state setting is that $\mu_t^\theta$ is now an infinite-dimensional object. In the discrete case, the population law is exactly represented by a point of a finite-dimensional simplex. No analogous exact representation exists for a general probability measure on $\R^d$. We therefore introduce a finite-dimensional approximation of the population law before applying the randomization argument.

Fix an integer $K\geq1$. We consider the family of $K$-component Gaussian mixtures
\begin{align}
    \Gc_K\coloneqq \left\{\sum_{k=1}^{K} w_k\Nc(m_l,\Sigma_k) : w\in\Delta_K^\circ,\;m_k\in\R^d,\;\Sigma_k\in\S_{++}^d \right\},
\end{align}
where $\Delta_K^\circ$ denotes the relative interior of the $K$-dimensional probability simplex and $\mathbb S_{++}^d$ the set of symmetric positive-definite $d\times d$ matrices.

To perturb elements of $\Gc_K$ while automatically preserving the constraints on the weights and covariance matrices, we work in unconstrained Euclidean coordinates. Let $q_K\coloneqq (K-1) + Kd + K\frac{d(d+1)}{2}$. For $z=\left(\beta,m_1,\ldots,m_K,s_1,\ldots,s_K \right)\in\R^{q_K}$, with $\beta\in\R^{K-1}$, define the mixture weights by the logistic parametrization
\begin{align}
    w_k(\beta) &= \frac{\exp(\beta_k)}{1 + \sum_{j=1}^{K-1}\exp(\beta_j)},\qquad k=1,\ldots,K-1,\\
    w_K(\beta) &= \frac{1}{1+\sum_{k=1}^{K-1} \exp(\beta_j)}.
\end{align}
For each $k$, the vector $s_k\in\R^{d(d+1)/2}$ parametrizes a lower-triangular matrix $L(s_k)$ whose diagonal entries are exponentials of the corresponding diagonal coordinates of $s_k$. We then set $\Sigma(s_k)\coloneqq L(s_k)L(s_k)^\top\in\S_{++}^d$. This defines a smooth decoder
\begin{align}
    \label{eq:gmm-decoder}
    \Gamma_K:\R^{q_K}\longrightarrow\Gc_K,\qquad \Gamma_K(z)\coloneqq\sum_{k=1}^{K} w_k(\beta)\Nc\left(m_k,\Sigma(s_k)\right).
\end{align}

The map $\Gamma_K$ transforms arbitrary Euclidean parameters into a valid probability measure. This plays the same role as the softmax map in the finite-state construction: perturbations can be performed in an unconstrained Euclidean space and subsequently mapped back to the space of admissible probability measures.

\paragraph{Finite-dimensional representation of the population flow.}

We introduce a measurable representation map $\Rc_K:\Pc_1(\R^d)\to\R^{q_K}$, and define
\begin{align}
    \label{eq:gmm-projected-law}
    z_{K,t}^\theta &\coloneqq \Rc_K(\mu_t^\theta),\qquad
    \bar\mu_{K,t}^\theta&\coloneqq\Gamma_K(z_{K,t}^\theta).
\end{align}
The measure $\bar\mu_{K,t}^\theta$ is the $K$-component Gaussian-mixture representation of the true population law $\mu_t^\theta$.

At this stage, $\Rc_K$ is kept abstract. It may correspond, for instance, to a population-level Gaussian-mixture fitting procedure. In the numerical implementation it will be replaced by an estimator based on particles sampled from $\mu_t^\theta$. The precise statistical procedure used to construct $\Rc_K$ is conceptually separate from the policy-gradient argument.

For a metric $d_{\Pc}$ on $\Pc_1(\R^d)$, define the representation error
\begin{align}
    \label{eq:gmm-representation-error}
    \delta_{K,t}^\theta\coloneqq d_\Pc \left( \mu_t^\theta,\bar\mu_{K,t}^\theta\right),\qquad \delta_K^\theta\coloneqq\max_{t\in\Tc}\delta_{K,t}^\theta.
\end{align}
Unlike the finite-state parametrization, this error is generally non-zero for fixed $K$. It constitutes an additional approximation error that is independent of the perturbation scale introduced below.

When $\mu_t^\theta\in\Gc_K$ and the representation map is exact, we have $\bar\mu_{K,t}^\theta=\mu_t^\theta$ with $\delta_{K,t}^\theta=0$. In particular, the Gaussian setting corresponds to the special case $K=1$. Thus the Gaussian linear--quadratic model studied in Chapter~\ref{chapter:experiments} is recovered as an exact finite-dimensional subcase of the present construction.

\begin{remark}[Identifiability of the mixture representation]
    The decoder $\Gamma_K$ is not injective because permutations of the mixture components represent the same probability measure. We therefore assume that $\Rc_K$ selects a deterministic representative, for example through a fixed component-ordering convention. For the differentiation arguments below, we additionally restrict attention to non-degenerate representations, with mixture weights bounded away from zero and covariance matrices uniformly positive definite. These conditions prevent label collisions and singular components along the population flow.
\end{remark}

\paragraph{Random perturbation of the Gaussian-mixture coordinates.}
Let $(U_t)_{t\in\Tc}$ be i.i.d.\ $\R^{q_K}$-valued random variables, independent of the state and policy noises, with common distribution $\nu$ admitting a density $\rho$ with respect to Lebesgue measure. For $\lambda>0$, define the randomized mixture parameter
\begin{align}
    \label{eq:gmm-random-parameter}
    Z_{K,t}^{\lambda,\theta}\coloneqq z_{K,t}^\theta + \lambda U_t,
\end{align}
and the corresponding randomized population law
\begin{align}
    \label{eq:gmm-random-law}
    M_{K,t}^{\lambda,\theta} \coloneqq \Gamma_K\left(Z_{K,t}^{\lambda,\theta}\right) = \Gamma_K\left( z_{K,t}^\theta+\lambda U_t\right).
\end{align}
By construction, $M_{K,t}^{\lambda,\theta}\in\Gc_K\subset\Pc_1(\R^d)$ almost surely. Moreover, $M_{K,t}^{0,\theta} = \bar\mu_{K,t}^\theta$.

It is important to distinguish the three objects $\mu_t^\theta$, $\bar\mu_{K,t}^\theta$, and $M_{K,t}^{\lambda,\theta}$. The first is the true population law of the unperturbed controlled process. The second is its deterministic finite-dimensional Gaussian-mixture representation. The third is the randomized population argument used to construct the perturbed control problem. In particular, $M_{K,t}^{\lambda,\theta}$ should not be identified with the law of the perturbed state process defined below.

We now define the perturbed controlled process by replacing the population argument in the policy and transition kernel with $M_{K,t}^{\lambda,\theta}$:
\begin{align}
    \label{eq:gmm-perturbed-process}
    \begin{cases}
        X_0^{K,\lambda,\theta}\sim\mu_0,\\
        \alpha_t^{K,\lambda,\theta}\sim\pi_t^\theta\left(\cdot\mid X_t^{K,\lambda,\theta},M_t^{K,\lambda,\theta}\right),\\
        X_{t+1}^{K,\lambda,\theta} \sim P_t\left(\cdot\mid X_t^{K,\lambda,\theta},M_{K,t}^{\lambda,\theta}, \alpha_t^{K,\lambda,\theta}\right),\qquad t=0,\ldots,T-1.
    \end{cases}
\end{align}
The corresponding objective is
\begin{align}
    \label{eq:gmm-perturbed-objective}
    J_K^\lambda(\theta)\coloneqq\E\left[\sum_{t=0}^{T-1}r_t\left(X_t^{K,\lambda,\theta}, M_{K,t}^{\lambda,\theta},\alpha_t^{K,\lambda,\theta}\right) + g\left(X_T^{K,\lambda,\theta}, M_{K,T}^{\lambda,\theta}\right) \right],
\end{align}
where the expectation is taken over the mixture perturbations, the state process, and the randomized actions.

For fixed $K$, sending $\lambda$ to zero removes the random perturbation but does not remove the finite-dimensional representation error. Consequently, the limit $\lambda\downarrow0$ should first be compared with the objective obtained from the deterministic represented population flow $(\bar\mu_{K,t}^\theta)_{t\in\Tc}$. Recovery of the original continuous-state MFC objective additionally requires the representation error $\delta_K^\theta$ to vanish as $K$ increases.

This gives two distinct approximation scales:
\begin{align}
    \label{eq:gmm-two-approximation-scales}
    \mu_t^\theta \xrightarrow[K\to\infty]{\text{representation}}\bar\mu_{K,t}^\theta,\qquad \bar\mu_{K,t}^\theta\xleftarrow[\lambda\downarrow0]{\text{perturbation}}M_{K,t}^{\lambda,\theta}.
\end{align}

The following elementary estimate makes this decomposition explicit.

\begin{assumption}[Regularity of the Gaussian-mixture decoder]
    \label{ass:gmm-decoder-regularity}
    Let $d_\Pc=W_1$. There exist $\lambda_0>0$ and $C_{\Gamma,K}<\infty$ such that, for all $\theta\in\Theta$, $t\in\Tc$, and $\lambda\in[0,\lambda_0]$,
    \begin{align}
        \label{eq:gmm-decoder-lipschitz}
        \E\left[W_1\left(\Gamma_K(z_{K,t}^\theta + \lambda U_t), \Gamma_K(z_{K,t}^\theta) \right) \right] \leq C_{\Gamma,K}\lambda.
    \end{align}
\end{assumption}

This assumption is a local regularity condition on the finite-dimensional chart. It is satisfied under uniform non-degeneracy and moment bounds on the mixture parameters together with suitable moment conditions on the perturbation distribution.

\begin{theorem}[Representation and perturbation estimate]
    \label{theorem:gmm-perturbation-estimate}
    Under Assumption~\ref{ass:gmm-decoder-regularity}, for every $\theta\in\Theta$, $t\in\Tc$, and $\lambda\in[0,\lambda_0]$,
    \begin{align}
        \label{eq:gmm-total-law-error}
        \E\left[W_1\left(M_{K,t}^{\lambda,\theta},\mu_t^\theta\right) \right] \leq \delta_{K,t}^\theta + C_{\Gamma,K}\lambda.
    \end{align}
    Consequently, if $\sup_{\theta\in\Theta}\max_{t\in\Tc}\delta_{K,t}^\theta\xrightarrow[K\to\infty]{}0$, then
    \begin{align}
        \sup_{\theta\in\Theta}\max_{t\in\Tc}\E\left[W_1\left(M_{K,t}^{\lambda,\theta}, \mu_t^\theta\right)\right]\xrightarrow[\substack{\lambda\downarrow0\\ K\to\infty}]{}0.
    \end{align}
\end{theorem}

\begin{proof}
    By the triangle inequality,
    \begin{align}
        W_1\left(M_{K,t}^{\lambda,\theta},\mu_t^\theta \right)
        &\leq W_1\left(M_{K,t}^{\lambda,\theta},\bar\mu_{K,t}^\theta\right) + W_1\left(\bar\mu_{K,t}^\theta,\mu_t^\theta\right)\\
        &= W_1\left(\Gamma_K(z_{K,t}^\theta+\lambda U_t), \Gamma_K(z_{K,t}^\theta)\right) + \delta_{K,t}^\theta.
    \end{align}
    Taking expectations and applying Assumption~\ref{ass:gmm-decoder-regularity} gives
    \begin{align}
        \E\left[ W_1\left(M_{K,t}^{\lambda,\theta},\mu_t^\theta\right)\right]\leq C_{\Gamma,K}\lambda + \delta_{K,t}^\theta.
    \end{align}
    The second statement follows immediately.
\end{proof}

The estimate \eqref{eq:gmm-total-law-error} is the continuous-state analogue of the simplex perturbation estimate obtained in the finite-state setting. There is, however, one additional term: $\delta_{K,t}^\theta$ measures the structural error caused by replacing an infinite-dimensional population law with a finite-dimensional Gaussian mixture. The subsequent convergence and policy-gradient analysis must therefore keep the limits in $K$ and $\lambda$ separate.

Finally, the use of the Euclidean coordinates $Z_{K,t}^{\lambda,\theta}$ has an additional advantage. We do not need to construct a dominating measure or a density directly on the infinite-dimensional space $\Pc(\R^d)$. The likelihood-ratio correction can instead be derived from the ordinary finite-dimensional density of $Z_{K,t}^{\lambda,\theta} = z_{K,t}^\theta + \lambda U_t$.

\subsection{Density and score of the randomized Gaussian-mixture parameters}
\label{subsec:gmm-density-score}

We now derive the density and score associated with the randomized finite-dimensional population representation introduced in Section~\ref{subsec:continuous-gmm-perturbation}. The main advantage of the Gaussian-mixture parametrization is that the randomization is performed directly in the Euclidean coordinate space $\R^{q_K}$ rather than on the infinite-dimensional space $\Pc(\R^d)$.

Recall that $Z_{K,t}^{\lambda,\theta} = z_{K,t}^{\theta} + \lambda U_t$ and $z_{K,t}^{\theta} = \Rc_K(\mu_t^\theta)$, where $(U_t)_{t\in\Tc}$ are i.i.d.\ $\R^{q_K}$-valued random variables with common density $\rho:\R^{q_K}\to\R_+$ with respect to Lebesgue measure. The randomized population argument is then $M_{K,t}^{\lambda,\theta} = \Gamma_K(Z_{K,t}^{\lambda,\theta})$.

The dependence of the distribution of $Z_{K,t}^{\lambda,\theta}$ on $\theta$ is entirely induced by the translation $z_{K,t}^{\theta}$. This yields an explicit finite-dimensional likelihood-ratio correction.

For $\lambda > 0$, define
\begin{align}
    \label{eq:gmm-coordinate-density}
    \varrho_{K,t}^{\lambda,\theta}(z) \coloneqq \frac{1}{\lambda^{q_K}}\rho\left(\frac{z-z_{K,t}^\theta}{\lambda}\right),\qquad z\in\R^{q_K}.
\end{align}

\begin{proposition}[Density of the randomized mixture coordinates]
    \label{prop:gmm-coordinate-density}
    Let $\lambda>0$ and suppose that $U_t$ admits density $\rho$ with respect to Lebesgue measure on $\R^{q_K}$. Then $Z_{K,t}^{\lambda,\theta}$ admits density $\varrho_{K,t}^{\lambda,\theta}$ given by \eqref{eq:gmm-coordinate-density}. In particular,
    \begin{align}
        \Lc(Z_{K,t}^{\lambda,\theta})(\d z) = \varrho_{K,t}^{\lambda,\theta}(z)\d z.
    \end{align}
\end{proposition}

\begin{proof}
    For fixed $(t,K,\lambda,\theta)$, consider the affine map
    \begin{align}
        \Phi_{K,t}^{\lambda,\theta}:\R^{q_K} &\longrightarrow\R^{q_K},\\
        u &\longmapsto z_{K,t}^\theta + \lambda u.
    \end{align}
    Its inverse is
    \begin{align}
        \left(\Phi_{K,t}^{\lambda,\theta}\right)^{-1}(z) = \frac{z-z_{K,t}^\theta}{\lambda},
    \end{align}
    and
    \begin{align}
        \left|\det D\left(\Phi_{K,t}^{\lambda,\theta}\right)^{-1}(z)\right| = \lambda^{-q_K}.
    \end{align}
    The change-of-variables formula therefore gives
    \begin{align}
        \varrho_{K,t}^{\lambda,\theta}(z) &= \rho\left(\left(\Phi_{K,t}^{\lambda,\theta}\right)^{-1}(z)\right)\left|\det D\left(\Phi_{K,t}^{\lambda,\theta}\right)^{-1}(z) \right| \\
        &= \frac{1}{\lambda^{q_K}}\rho\left(\frac{z-z_{K,t}^{\theta}}{\lambda} \right).
    \end{align}
\end{proof}

\begin{remark}[Why we work with $Z_{K,t}^{\lambda,\theta}$]
    \label{rem:no-density-on-gmm-space}
    We do not require the random measure $M_{K,t}^{\lambda,\theta} = \Gamma_K(Z_{K,t}^{\lambda,\theta})$ itself to admit a density with respect to a reference measure on $\Pc(\R^d)$ or even on $\Gc_K$.
    
    Indeed, because the Gaussian-mixture decoder $\Gamma_K$ is non-injective under permutations of the mixture components, constructing such a density directly on $\Gc_K$ would introduce unnecessary technical complications. Instead, we retain the Euclidean random variable $Z_{K,t}^{\lambda,\theta}$ as part of the augmented trajectory and evaluate the policy, transition kernel, and rewards through the decoded measure $\Gamma_K(Z_{K,t}^{\lambda,\theta})$.
    
    Thus, all likelihood-ratio calculations can be carried out with respect to ordinary Lebesgue measure on $\R^{q_K}$.
\end{remark}

Let the policy parameter satisfy $\theta\in\Theta\subset\R^{d_\theta}$, and suppose that $\theta\longmapsto z_{K,t}^\theta\in\R^{q_K}$ is differentiable. We denote its Jacobian by $D_\theta z_{K,t}^\theta\in\R^{q_K\times d_\theta}$.

We now compute the score of the family of densities $\left(\varrho_{K,t}^{\lambda,\theta}\right)_{\theta\in\Theta}$.

\begin{proposition}[Score of the randomized mixture coordinates]
    \label{prop:gmm-coordinate-score}
    Suppose that $\rho$ is strictly positive and continuously differentiable on $\R^{q_K}$ and that $\theta\mapsto z_{K,t}^{\theta}$ is differentiable. Then, for every $z\in\R^{q_K}$,
    \begin{align}
        \label{eq:gmm-coordinate-score-general}
        \nabla_\theta\log\varrho_{K,t}^{\lambda,\theta}(z) = -\frac{1}{\lambda}\left(D_\theta z_{K,t}^\theta\right)^\top\nabla\log\rho\left(\frac{z-z_{K,t}^\theta}{\lambda}\right).
    \end{align}
    Consequently, evaluated at the sampled randomized coordinate $Z_{K,t}^{\lambda,\theta} = z_{K,t}^\theta + \lambda U_t$, we obtain
    \begin{align}
        \label{eq:gmm-coordinate-score-sampled}
        \nabla_\theta\log\varrho_{K,t}^{\lambda,\theta}\left(Z_{K,t}^{\lambda,\theta}\right) = -\frac{1}{\lambda}\left(D_\theta z_{K,t}^\theta\right)^\top\nabla\log\rho(U_t).
    \end{align}
\end{proposition}

\begin{proof}
    From Proposition~\ref{prop:gmm-coordinate-density},
    \begin{align}
        \log\varrho_{K,t}^{\lambda,\theta}(z) = -q_K\log\lambda + \log\rho\left(\frac{z-z_{K,t}^\theta}{\lambda}\right).
    \end{align}
    For fixed $z$, the first term does not depend on $\theta$. By the chain rule,
    \begin{align}
        \nabla_\theta\log\varrho_{K,t}^{\lambda,\theta}(z)
        &= \left[D_\theta\left(\frac{z-z_{K,t}^{\theta}}{\lambda}\right)\right]^\top \nabla\log\rho\left(\frac{z-z_{K,t}^\theta}{\lambda}\right)\\
        &= -\frac{1}{\lambda}\left(D_\theta z_{K,t}^{\theta}\right)^\top \nabla\log\rho\left(\frac{z-z_{K,t}^\theta}{\lambda}\right).
    \end{align}
    Evaluating at $z=Z_{K,t}^{\lambda,\theta}=z_{K,t}^{\theta}+\lambda U_t$ gives \eqref{eq:gmm-coordinate-score-sampled}.
\end{proof}

The score therefore depends on the unknown mean-field dynamics only through the population-flow sensitivity $D_\theta z_{K,t}^{\theta} = D_\theta \Rc_K(\mu_t^\theta)$. This is the continuous-state counterpart of the gradient of the logits in the finite-state construction. A model-free estimator of $D_\theta z_{K,t}^{\theta}$ will be developed subsequently.

A particularly convenient choice is to use Gaussian perturbations. Let $U_t\sim\Nc(0,\Xi)$ for $\Xi\in\S_{++}^{q_K}$, independently across $t$. Then
\begin{align}
    \rho(u)=\frac{1}{(2\pi)^{q_K/2} \det(\Xi)^{1/2}}\exp\left(-frac12 u^\top\Xi^{-1}u\right),
\end{align}
and therefore $\nabla\log\rho(u) = -\Xi^{-1}u$.

Corollary~\ref{prop:gmm-coordinate-score} then becomes particularly simple.

\begin{corollary}[Gaussian perturbation score]
    \label{cor:gmm-gaussian-score}
    Suppose that $U_t\sim\Nc(0,\Xi)$. Then
    \begin{align}
        \label{eq:gmm-gaussian-score}
        \nabla_\theta\log\varrho_{K,t}^{\lambda,\theta}\left(Z_{K,t}^{\lambda,\theta}\right) = \frac1\lambda \left(D_\theta z_{K,t}^\theta\right)^\top \Xi^{-1}U_t.
    \end{align}
    In particular, for $U_t\sim\Nc(0,I_{q_K})$, we have
    \begin{align}
        \label{eq:gmm-isotropic-score}
        \nabla_\theta\log\varrho_{K,t}^{\lambda,\theta}\left(Z_{K,t}^{\lambda,\theta}\right) = \frac1\lambda\left(D_\theta z_{K,t}^\theta\right)^\top U_t.
    \end{align}
\end{corollary}

Equation~\eqref{eq:gmm-isotropic-score} is the direct continuous-state analogue of the perturbation-score term obtained in the finite-state setting. The singular factor $\lambda^{-1}$ is also the source of the familiar bias--variance trade-off: reducing $\lambda$ decreases the perturbation bias, while simultaneously amplifying the variance of likelihood-ratio estimators.

\begin{remark}[Anisotropic perturbations]
    \label{rem:gmm-anisotropic-noise}
    Although the isotropic Gaussian choice is the simplest theoretically, the coordinates of the Gaussian-mixture representation have different geometries and scales. In particular, mixture logits, component means, and covariance coordinates need not respond similarly to perturbations.
    
    It can therefore be advantageous in practice to use a block-diagonal covariance $\Xi=\operatorname{diag}\left(\Xi_w,\Xi_m,\Xi_\Sigma\right)$, corresponding respectively to mixture-weight, mean, and covariance coordinates. The score formula \eqref{eq:gmm-gaussian-score} remains explicit in this case.
\end{remark}

The usual likelihood-ratio identities hold for the perturbation score. Define $S_{K,t}^{\lambda,\theta} \coloneqq \nabla_\theta \log \varrho_{K,t}^{\lambda,\theta} \left( Z_{K,t}^{\lambda,\theta} \right)$. Then, under the usual integrability assumptions, $\E \left[ S_{K,t}^{\lambda,\theta} \right] = 0$. Indeed,
\begin{align}
    \E \left[ S_{K,t}^{\lambda,\theta} \right] &= \int_{\R^{q_K}} \nabla_\theta \log \varrho_{K,t}^{\lambda,\theta}(z) \, \varrho_{K,t}^{\lambda,\theta}(z) \,dz \\ &= \int_{\R^{q_K}} \nabla_\theta \varrho_{K,t}^{\lambda,\theta}(z) \,dz \\ &= \nabla_\theta \int_{\R^{q_K}} \varrho_{K,t}^{\lambda,\theta}(z) \,dz = 0.
\end{align}

More generally, if $H:\R^{q_K}\to\R$ is sufficiently integrable and does not depend explicitly on $\theta$, then
\begin{align}
    \nabla_\theta \E \left[ H \left( Z_{K,t}^{\lambda,\theta} \right) \right] = \E \left[ H \left( Z_{K,t}^{\lambda,\theta} \right) S_{K,t}^{\lambda,\theta} \right].
    \label{eq:gmm-likelihood-ratio-identity}
\end{align}
This identity is the mechanism by which derivatives of the population-dependent reward, policy, and transition kernel will be replaced by a model-free score term.

\paragraph{Augmented formulation.}
For later use, it is convenient to regard the randomized coordinate $Z_{K,t}^{\lambda,\theta}$ as an explicit component of the perturbed trajectory. Define
\begin{align}
    \widetilde r_{K,t}(x,a,z) &\coloneqq r_t \left( x, \Gamma_K(z), a \right), \\ \widetilde g_K(x,z) &\coloneqq g \left( x, \Gamma_K(z) \right), \\ \widetilde\pi_{K,t}^{\theta}(da\mid x,z) &\coloneqq \pi_t^\theta \left( da \mid x, \Gamma_K(z) \right), \\ \widetilde P_{K,t}(dx'\mid x,a,z) &\coloneqq P_t \left( dx' \mid x, a, \Gamma_K(z) \right).
    \label{eq:gmm-augmented-coefficients}
\end{align}

The perturbed objective may therefore be written entirely on the finite-dimensional augmented space $\R^d \times \R^{q_K} \times \Ac$. The corresponding trajectory law takes the form
\begin{align}
    \mathbb P_{K,\lambda}^{\theta} &(dx_0,dz_0,da_0,\ldots, dx_{T-1},dz_{T-1},da_{T-1},dx_T,dz_T) \nonumber\\ &= \mu_0(dx_0) \prod_{t=0}^{T-1} \Big[ \varrho_{K,t}^{\lambda,\theta}(z_t)\,dz_t\, \widetilde\pi_{K,t}^{\theta}(da_t\mid x_t,z_t)\, \widetilde P_{K,t}(dx_{t+1}\mid x_t,a_t,z_t) \Big] \, \varrho_{K,T}^{\lambda,\theta}(z_T)\,dz_T.
    \label{eq:gmm-augmented-trajectory-law}
\end{align}

This representation makes the key structural point explicit: after augmenting the trajectory by the randomized finite-dimensional coordinates $(Z_{K,t}^{\lambda,\theta})_{t\in\Tc}$, the dependence of the trajectory density on $\theta$ occurs through only two families of terms, $\widetilde\pi_{K,t}^{\theta}$ and $\varrho_{K,t}^{\lambda,\theta}$. In particular, the model transition kernel does not need to be differentiated. This observation will yield the model-free policy-gradient representation of the perturbed continuous-state control problem.

\subsection{Convergence of the perturbed objective}
\label{subsec:gmm-objective-convergence}

We now show that the Gaussian-mixture perturbed control problem approximates the original mean-field control problem. There are two distinct sources of approximation. First, for a fixed number $K$ of Gaussian components, the random perturbation $M_{K,t}^{\lambda,\theta} = \Gamma_K \left( z_{K,t}^\theta+\lambda U_t \right)$ approaches the deterministic Gaussian-mixture representation $\bar\mu_{K,t}^\theta$ as $\lambda\downarrow0$. Second, even when $\lambda=0$, the represented law $\bar\mu_{K,t}^\theta$ need not coincide with the true population law $\mu_t^\theta$.

Accordingly, we decompose
\begin{align}
    \label{eq:gmm-objective-error-decomposition}
    J_K^\lambda(\theta) - J(\theta) = \left( J_K^\lambda(\theta)-J_K^0(\theta)\right) + \left(J_K^0(\theta)-J(\theta)\right),
\end{align}
where $J_K^0$ denotes the control problem obtained by replacing the population argument by the deterministic Gaussian-mixture representation $\bar\mu_{K,t}^\theta$.

The first term in \eqref{eq:gmm-objective-error-decomposition} is the perturbation error, whereas the second is the finite-dimensional representation error.

\paragraph{Averaged controlled kernel and reward.}

For $\theta\in\Theta$, define the averaged controlled transition kernel
\begin{align}
    \Kc_t^\theta\left( dx' \mid x,m \right) \coloneqq \int_{\Ac} P_t \left( dx' \mid x,a,m \right) \pi_t^\theta \left( da \mid x,m \right),
    \label{eq:gmm-averaged-kernel}
\end{align}
and the averaged running reward
\begin{align}
    \Rc_t^\theta(x,m) \coloneqq \int_{\Ac} r_t(x,a,m) \pi_t^\theta \left( da \mid x,m \right).
    \label{eq:gmm-averaged-reward}
\end{align}

The law $\mu_t^\theta$ of the original controlled process therefore satisfies
\begin{align}
    \mu_{t+1}^\theta = \int_{\R^d} \Kc_t^\theta \left( \cdot \mid x,\mu_t^\theta \right) \mu_t^\theta(dx).
    \label{eq:gmm-original-flow}
\end{align}

We denote by $\nu_{K,t}^{\lambda,\theta} \coloneqq \Lc \left( X_t^{K,\lambda,\theta} \right)$ the law of the state process driven by the randomized population arguments $(M_{K,t}^{\lambda,\theta})_{t\in\Tc}$.

Since $U_t$ is independent of the state and action randomness up to time $t$, the randomized measure $M_{K,t}^{\lambda,\theta}$ is independent of $X_t^{K,\lambda,\theta}$. Consequently,
\begin{align}
    \nu_{K,t+1}^{\lambda,\theta} = \E \left[ \Kc_t^\theta \left( \cdot \mid X_t^{K,\lambda,\theta}, M_{K,t}^{\lambda,\theta} \right) \right].
    \label{eq:gmm-perturbed-flow}
\end{align}

When $\lambda=0$, we similarly denote $\nu_{K,t}^{0,\theta} = \Lc \left( X_t^{K,0,\theta} \right)$, where the population argument is the deterministic measure $\bar\mu_{K,t}^\theta$.

Notice that, in general, $\nu_{K,t}^{0,\theta}\neq\bar\mu_{K,t}^\theta$. The latter is an approximation of the original population law $\mu_t^\theta$, whereas the former is the law of the state process obtained after using this approximation as an exogenous population argument.

\begin{assumption}[Wasserstein stability of the controlled dynamics]
    \label{ass:gmm-objective-stability}
    There exist constants $L_{\Kc,x}, L_{\Kc,\mu}, L_{\Rc,x}, L_{\Rc,\mu}, L_{g,x}, L_{g,\mu} \geq 0$ such that, uniformly in $\theta\in\Theta$ and $t=0,\ldots,T-1$,
    \begin{align}
        W_1\left(\Kc_t^\theta(\cdot\mid x,m), \Kc_t^\theta(\cdot\mid x',m') \right) &\leq L_{\Kc,x}|x-x'| + L_{\Kc,\mu}W_1(m,m'), \label{eq:gmm-kernel-lipschitz}\\
        \left| \Rc_t^\theta(x,m) - \Rc_t^\theta(x',m') \right| &\leq L_{\Rc,x}|x-x'| + L_{\Rc,\mu}W_1(m,m'), \label{eq:gmm-reward-lipschitz}\\
        \left| g(x,m) - g(x',m') \right| &\leq L_{g,x}|x-x'| + L_{g,\mu}W_1(m,m'). \label{eq:gmm-terminal-lipschitz}
    \end{align}
    We further assume the moment conditions required for all the Wasserstein distances and expectations below to be finite.
\end{assumption}

\begin{remark}[Primitive sufficient conditions]
    Assumption~\ref{ass:gmm-objective-stability} is formulated directly in terms of the averaged controlled coefficients because these are the quantities entering the population recursion. It can be verified from primitive Lipschitz assumptions on $P_t$, $\pi_t^\theta$, and $r_t$.
    
    For instance, sufficient conditions for \eqref{eq:gmm-kernel-lipschitz} can be obtained by combining Wasserstein regularity of $(x,m) \longmapsto P_t(\cdot\mid x,a,m)$ uniformly in $a$, with suitable regularity of the randomized policy $(x,m) \longmapsto \pi_t^\theta(\cdot\mid x,m)$. We work directly with $\Kc_t^\theta$ and $\Rc_t^\theta$ in order to keep the stability argument independent of the particular action-space parametrization.
\end{remark}

We first record a general stability estimate that will be used both for the perturbation error and for the Gaussian-mixture representation error.

Let $(m_t)_{t\in\Tc}$ and $(\widetilde m_t)_{t\in\Tc}$ be two possibly random population-argument flows, independent of the corresponding state noises at the current time. Let $(Y_t)$ and $(\widetilde Y_t)$ be controlled processes driven by the same policy parameter $\theta$ and by $m_t$ and $\widetilde m_t$, respectively. Denote $\eta_t=\Lc(Y_t), \qquad \widetilde\eta_t=\Lc(\widetilde Y_t)$, and set $d_t\coloneqq\E\left[W_1(m_t,\widetilde m_t)\right]$.

\begin{lemma}[Stability of the state marginal]
    \label{lem:gmm-state-stability}
    Under Assumption~\ref{ass:gmm-objective-stability}, if $\eta_0=\widetilde\eta_0=\mu_0$, then
    \begin{align}
        W_1(\eta_t,\widetilde\eta_t) \leq L_{\Kc,\mu} \sum_{s=0}^{t-1} L_{\Kc,x}^{\,t-1-s} d_s, \qquad t=1,\ldots,T.
        \label{eq:gmm-state-stability}
    \end{align}
\end{lemma}

\begin{proof}
    We proceed recursively. Since both processes have the same initial law, $W_1(\eta_0,\widetilde\eta_0)=0$.

    Let $e_t\coloneqq W_1(\eta_t,\widetilde\eta_t)$. Choose a coupling $(Y_t,\widetilde Y_t)$ such that $\E\left[|Y_t-\widetilde Y_t|\right]=e_t$. Conditionally on $(Y_t,\widetilde Y_t,m_t,\widetilde m_t)$, couple the next states through an optimal coupling of the two controlled transition kernels. By \eqref{eq:gmm-kernel-lipschitz},
    \begin{align}
        \label{eq:gmm-state-recursion}
        e_{t+1}&\leq\E\left[W_1\left(\Kc_t^\theta(\cdot\mid Y_t,m_t),\Kc_t^\theta(\cdot\mid \widetilde Y_t,\widetilde m_t)\right)\right]\\
        &\leq L_{\Kc,x}\E\left[|Y_t-\widetilde Y_t|\right] +L_{\Kc,\mu}\E\left[W_1(m_t,\widetilde m_t)\right]\\
        &= L_{\Kc,x}e_t + L_{\Kc,\mu}d_t.
    \end{align}
    Iterating \eqref{eq:gmm-state-recursion} from $e_0=0$ gives \eqref{eq:gmm-state-stability}.
\end{proof}

For convenience, define
\begin{align}
    H_t(L) \coloneqq
    \begin{cases}
        0, &t=0, \\ \displaystyle \sum_{j=0}^{t-1}L^j, &t\geq1.
    \end{cases}
    \label{eq:gmm-geometric-factor}
\end{align}
Thus, if $d_s\leq d$ uniformly over $s$, Lemma~\ref{lem:gmm-state-stability} gives $W_1(\eta_t,\widetilde\eta_t) \leq L_{\Kc,\mu} H_t(L_{\Kc,x})d$.

Recall that $M_{K,t}^{\lambda,\theta} = \Gamma_K \left( z_{K,t}^\theta+\lambda U_t \right)$ and $\bar\mu_{K,t}^\theta = \Gamma_K(z_{K,t}^\theta)$. Under Assumption~\ref{ass:gmm-decoder-regularity},
\begin{align}
    \E \left[ W_1 \left( M_{K,t}^{\lambda,\theta}, \bar\mu_{K,t}^\theta \right) \right] \leq C_{\Gamma,K}\lambda.
    \label{eq:gmm-perturbation-distance-recall}
\end{align}

Applying Lemma~\ref{lem:gmm-state-stability} with $m_t=M_{K,t}^{\lambda,\theta}$, and $\widetilde m_t=\bar\mu_{K,t}^\theta$, gives
\begin{align}
    W_1 \left( \nu_{K,t}^{\lambda,\theta}, \nu_{K,t}^{0,\theta} \right) \leq L_{\Kc,\mu} H_t(L_{\Kc,x}) C_{\Gamma,K}\lambda.
    \label{eq:gmm-perturbed-state-to-unperturbed-gmm}
\end{align}

We can now control the objective.
\begin{proposition}[Perturbation error]
    \label{prop:gmm-perturbation-objective}
    Under Assumptions~\ref{ass:gmm-decoder-regularity} and \ref{ass:gmm-objective-stability}, for every fixed $K$, $\theta\in\Theta$, and $\lambda\in[0,\lambda_0]$,
    \begin{align}
        \left| J_K^\lambda(\theta) - J_K^0(\theta) \right| \leq C_T C_{\Gamma,K}\lambda, \label{eq:gmm-perturbation-objective-bound}
    \end{align}
    where
    \begin{align}
        C_T \coloneqq \sum_{t=0}^{T-1} \left[ L_{\Rc,x} L_{\Kc,\mu} H_t(L_{\Kc,x}) + L_{\Rc,\mu} \right] + L_{g,x} L_{\Kc,\mu} H_T(L_{\Kc,x}) + L_{g,\mu}.
        \label{eq:gmm-objective-stability-constant}
    \end{align}
    Consequently, for fixed $K$,
    \begin{align}
        J_K^\lambda(\theta) \xrightarrow[\lambda\downarrow0]{} J_K^0(\theta).
    \end{align}
\end{proposition}

\begin{proof}
    For each running reward, using \eqref{eq:gmm-reward-lipschitz},
    \begin{align}
        & \left| \E \left[ \Rc_t^\theta \left( X_t^{K,\lambda,\theta}, M_{K,t}^{\lambda,\theta} \right) \right] - \E \left[ \Rc_t^\theta \left( X_t^{K,0,\theta}, \bar\mu_{K,t}^\theta \right) \right] \right| \nonumber\\
        &\qquad\leq L_{\Rc,x} W_1 \left( \nu_{K,t}^{\lambda,\theta}, \nu_{K,t}^{0,\theta} \right) + L_{\Rc,\mu} \E \left[ W_1 \left( M_{K,t}^{\lambda,\theta}, \bar\mu_{K,t}^\theta \right) \right].
    \end{align}
    Using \eqref{eq:gmm-perturbed-state-to-unperturbed-gmm} and \eqref{eq:gmm-perturbation-distance-recall}, this is bounded by
    \begin{align}
        \left[ L_{\Rc,x} L_{\Kc,\mu} H_t(L_{\Kc,x}) + L_{\Rc,\mu} \right] C_{\Gamma,K}\lambda.
    \end{align}
    
    Similarly, the terminal reward satisfies
    \begin{align}
        & \left| \E \left[ g \left( X_T^{K,\lambda,\theta}, M_{K,T}^{\lambda,\theta} \right) \right] - \E \left[ g \left( X_T^{K,0,\theta}, \bar\mu_{K,T}^\theta \right) \right] \right| \\
        &\qquad\leq \left[ L_{g,x} L_{\Kc,\mu} H_T(L_{\Kc,x}) + L_{g,\mu} \right] C_{\Gamma,K}\lambda.
    \end{align}
    Summing over time gives \eqref{eq:gmm-perturbation-objective-bound}.
\end{proof}

We now compare $J_K^0$ with the original objective $J$. Recall that $\delta_{K,t}^\theta = W_1 \left( \bar\mu_{K,t}^\theta, \mu_t^\theta \right)$, and $\delta_K^\theta = \max_{t\in\Tc} \delta_{K,t}^\theta$. Applying Lemma~\ref{lem:gmm-state-stability} with $m_t=\bar\mu_{K,t}^\theta$ and $\widetilde m_t=\mu_t^\theta$, we obtain
\begin{align}
    W_1 \left( \nu_{K,t}^{0,\theta}, \mu_t^\theta \right) \leq L_{\Kc,\mu} H_t(L_{\Kc,x}) \delta_K^\theta.
    \label{eq:gmm-representation-state-error}
\end{align}

\begin{proposition}[Representation error]
    \label{prop:gmm-representation-objective}
    Under Assumption~\ref{ass:gmm-objective-stability},
    \begin{align}
        \left| J_K^0(\theta) - J(\theta) \right| \leq C_T\delta_K^\theta,
        \label{eq:gmm-representation-objective-bound}
    \end{align} where $C_T$ is given by \eqref{eq:gmm-objective-stability-constant}.
\end{proposition}

\begin{proof}
    The proof is identical to that of Proposition~\ref{prop:gmm-perturbation-objective}, replacing $M_{K,t}^{\lambda,\theta}$ by $\bar\mu_{K,t}^\theta$ and $\bar\mu_{K,t}^\theta$ by $\mu_t^\theta$.

    In particular,
    \begin{align}
        & \left| \E \left[ \Rc_t^\theta \left( X_t^{K,0,\theta}, \bar\mu_{K,t}^\theta \right) \right] - \E \left[ \Rc_t^\theta \left( X_t^\theta, \mu_t^\theta \right) \right] \right| \\
        &\qquad\leq \left[ L_{\Rc,x} L_{\Kc,\mu} H_t(L_{\Kc,x}) + L_{\Rc,\mu} \right] \delta_K^\theta.
    \end{align}
    The analogous bound for the terminal reward and summation over time yield \eqref{eq:gmm-representation-objective-bound}.
\end{proof}

Combining the two preceding propositions gives the main result of this subsection.

\begin{theorem}[Convergence of the Gaussian-mixture perturbed objective]
    \label{thm:gmm-objective-convergence}
    Under Assumptions~\ref{ass:gmm-decoder-regularity} and \ref{ass:gmm-objective-stability}, for every $\theta\in\Theta$,
    \begin{align}
        \left| J_K^\lambda(\theta) - J(\theta) \right| \leq C_T \left( C_{\Gamma,K}\lambda + \delta_K^\theta \right).
        \label{eq:gmm-total-objective-bound}
    \end{align}
    In particular, for fixed $K$,
    \begin{align}
        \limsup_{\lambda\downarrow0} \left| J_K^\lambda(\theta) - J(\theta) \right| \leq C_T\delta_K^\theta.
        \label{eq:gmm-fixed-k-bias}
    \end{align}
    Furthermore, if $(K_n,\lambda_n)$ is a sequence satisfying
    \begin{align}
        \delta_{K_n}^\theta \longrightarrow0, \qquad C_{\Gamma,K_n}\lambda_n \longrightarrow0,
        \label{eq:gmm-joint-limit-condition}
    \end{align}
    then
    \begin{align}
        J_{K_n}^{\lambda_n}(\theta) \longrightarrow J(\theta).
        \label{eq:gmm-objective-joint-convergence}
    \end{align}
\end{theorem}

\begin{remark}[Order of the two limits]
    For fixed $K$, taking $\lambda\downarrow0$ only removes the artificial randomization. It does not recover the original MFC problem unless the true population flow belongs to the $K$-component Gaussian-mixture family.
    
    Thus, $J_K^\lambda \xrightarrow[\lambda\downarrow0]{} J_K^0$ while $J_K^0 \xrightarrow[K\to\infty]{} J$ requires consistency of the Gaussian-mixture representation.
    
    If the decoder constants $C_{\Gamma,K}$ are uniformly bounded in $K$, then any sequence satisfying $\lambda_K\to0$ and $\delta_K^\theta\to0$ is sufficient. Without such a uniform bound, the correct requirement is $C_{\Gamma,K}\lambda_K\to0$.
\end{remark}

\begin{remark}[Exact finite-dimensional closure]
    If the population flow belongs exactly to $\Gc_K$, so that $\mu_t^\theta = \bar\mu_{K,t}^\theta \qquad \text{for all }t$, then $\delta_K^\theta=0$ and 
    \begin{align}
        \left| J_K^\lambda(\theta) - J(\theta) \right| \leq C_T C_{\Gamma,K}\lambda.
    \end{align}
    This includes the Gaussian linear--quadratic setting with $K=1$ as a special case.
\end{remark}

For policy optimization, pointwise convergence in $\theta$ is not sufficient by itself. Uniform convergence over the policy class allows us to transfer optimality properties from the perturbed problems to the original MFC problem.

Define $\Delta_K \coloneqq \sup_{\theta\in\Theta} \delta_K^\theta$.

\begin{corollary}[Uniform convergence]
    \label{cor:gmm-uniform-objective-convergence} Suppose that the constants in Assumptions \ref{ass:gmm-decoder-regularity} and \ref{ass:gmm-objective-stability} are uniform in $\theta$. Then
    \begin{align}
        \sup_{\theta\in\Theta} \left| J_K^\lambda(\theta) - J(\theta) \right| \leq C_T \left( C_{\Gamma,K}\lambda + \Delta_K \right).
        \label{eq:gmm-uniform-objective-bound}
    \end{align}

    Consequently, if $\Delta_{K_n}\to0$ and $C_{\Gamma,K_n}\lambda_n\to0$, then
    \begin{align}
        \sup_{\theta\in\Theta} \left| J_{K_n}^{\lambda_n}(\theta) - J(\theta) \right| \longrightarrow0.
        \label{eq:gmm-uniform-objective-convergence}
    \end{align}
\end{corollary}

As a direct consequence, the optimal values converge.

\begin{corollary}[Convergence of optimal values and policies]
    \label{cor:gmm-optimal-policy-convergence}
    Assume the hypotheses of Corollary~\ref{cor:gmm-uniform-objective-convergence} and suppose that $\Theta$ is compact and that $J:\Theta\to\R$ is continuous. Let
    \begin{align}
        V &\coloneqq \inf_{\theta\in\Theta} J(\theta), \\
        V_K^\lambda &\coloneqq \inf_{\theta\in\Theta} J_K^\lambda(\theta).
    \end{align}
    Then
    \begin{align}
        \left| V_K^\lambda-V \right| \leq C_T \left( C_{\Gamma,K}\lambda+\Delta_K \right).
        \label{eq:gmm-optimal-value-bound}
    \end{align}

    Hence, for any sequence $(K_n,\lambda_n)$ satisfying $\Delta_{K_n}\to0$ and $C_{\Gamma,K_n}\lambda_n\to0$, we have $V_{K,n}^{\lambda_n}\to V$.

    Moreover, let $\theta_n^\star\in\argmin_{\theta\in\Theta}J_{K_n}^{\lambda_n}(\theta)$. Every accumulation point of $(\theta_n^\star)_{n\geq1}$ belongs to $\argmin_{\theta\in\Theta}J(\theta)$. In particular, if $J$ admits a unique minimizer $\theta^\star$, then
    \begin{align}
        \theta_n^\star \longrightarrow \theta^\star.
    \end{align}
\end{corollary}

\begin{proof}
    The optimal-value estimate follows from the elementary inequality $\left| \inf_\theta f(\theta) - \inf_\theta g(\theta) \right| \leq \sup_\theta |f(\theta)-g(\theta)|$.

    Let now $\theta_{n_j}^\star\to\bar\theta$ be an arbitrary convergent subsequence, which exists by compactness of $\Theta$. For every $\theta\in\Theta$, $J_{K_{n_j}}^{\lambda_{n_j}} (\theta_{n_j}^\star) \leq J_{K_{n_j}}^{\lambda_{n_j}} (\theta)$. By uniform convergence, $J(\theta_{n_j}^\star) - J_{K_{n_j}}^{\lambda_{n_j}} (\theta_{n_j}^\star) \longrightarrow0$ uniformly, and similarly for the right-hand side. Taking $j\to\infty$ and using continuity of $J$ gives $J(\bar\theta)\leq J(\theta)$ for $\theta\in\Theta$. Therefore $\bar\theta$ is an optimizer of the original problem. If the optimizer is unique, every convergent subsequence has the same limit $\theta^\star$, and compactness therefore implies $\theta_n^\star\to\theta^\star$.
\end{proof}

Theorem~\ref{thm:gmm-objective-convergence} establishes convergence at the objective level. For policy-gradient learning, however, we require the stronger property $\nabla_\theta J_K^\lambda(\theta) \longrightarrow \nabla_\theta J(\theta)$. Objective convergence alone does not imply convergence of gradients. 

We have so far introduced the Gaussian-mixture representation and perturbation of the population law, derived the density and score of the randomized GMM parameters, and proved the convergence of the perturbed objective.

The continuous-state theory is therefore still incomplete at this stage. The remaining work includes proving the gradient-level convergence, obtaining a policy-gradient formula for the perturbed problem, estimating the population-flow sensitivity, providing a Monte Carlo estimator of the perturbed policy gradient, and finally obtaining a bias and mean-square error decomposition. Additional benchmarks, algorithmic variants and numerical methods are also being considered, and may be incorporated during the final months of the internship.
